% route_t.tex -- the Route T note (content proposed for Phase C3 of ugjs/PLAN.md).
% Drafted 23 November 2026 by the drafter "route-t" of the working council (an AI agent).
% Audited by two independent agents on 23 November 2026, re-checked by both (23/24 November) and
% confirmed by one on 24 November; revised for their findings.
% Not integrated into the record.  Built on its own.
\documentclass[11pt,a4paper]{article}
\usepackage[margin=2.8cm]{geometry}
\usepackage{amsmath,amssymb,amsthm,mathrsfs}
\usepackage[T1]{fontenc}
\usepackage{microtype}
\usepackage{booktabs}
\usepackage{array}
\setlength{\emergencystretch}{3em}
\usepackage[colorlinks=true,linkcolor=blue,citecolor=blue,urlcolor=blue]{hyperref}

% --- macros of part_i/main.tex, same names and meanings (copied from ugjs/preamble.tex) ---
\newcommand{\R}{\mathbb{R}}
\newcommand{\Z}{\mathbb{Z}}
\newcommand{\C}{\mathbb{C}}
\newcommand{\E}{\mathbb{E}}
\newcommand{\A}{\mathfrak{A}}
\newcommand{\G}{\mathcal{G}}
\newcommand{\CE}{\mathcal{C}_{\mathrm E}}
\newcommand{\CH}{\mathcal{C}_{\mathrm H}}
\newcommand{\HH}{\mathcal{H}}
\newcommand{\Om}{\Omega}
\newcommand{\ip}[2]{\langle #1,#2\rangle}
\DeclareMathOperator{\supp}{supp}
\DeclareMathOperator{\dist}{dist}
\DeclareMathOperator{\spec}{spec}
\DeclareMathOperator{\Ran}{Ran}
\DeclareMathOperator{\Cov}{Cov}
\newcommand{\GJS}[1]{\lVert #1\rVert_{\mathrm{GJS}}}
\newcommand{\ev}[2]{\langle #1\rangle_{#2}}
\newcommand{\Uc}{\mathcal U}
\newcommand{\dmu}{d\mu_0}
\newcommand{\dphi}{d\phi_0}
% --- macros of this note ---
\newcommand{\Vb}{\overline V}                 % the time-integrated interaction of Lemma T.13
\newcommand{\Mt}{\widetilde M}                % the resolvent-weighted moment of Definition T.1
\newcommand{\Ft}{\widetilde F}
\newcommand{\Gt}{\widetilde G}
\newcommand{\mt}{\widetilde{\mathfrak m}}
\newcommand{\Psit}{\widetilde\Psi}

\renewcommand{\theequation}{t\arabic{equation}}
\theoremstyle{plain}
\newtheorem{tlemma}{Lemma}
\renewcommand{\thetlemma}{T.\arabic{tlemma}}
\newtheorem{ttheorem}[tlemma]{Theorem}
\newtheorem{tprop}[tlemma]{Proposition}
\theoremstyle{definition}
\newtheorem{tdef}[tlemma]{Definition}
\newtheorem{trem}[tlemma]{Remark}

\title{Route T: items (1)--(4) of Theorem 12.2 without the transverse sharp-time case\\[4pt]
\large A note for Phase C3 of \texttt{ugjs/PLAN.md}}
\author{Written by an AI agent (Claude), commissioned by the author\\
\small and supervised by the council's orchestrator (an AI agent), for Phase C3 of the project;\\
\small revised after two independent audits and their re-checks}
\date{23 November 2026}

\begin{document}
\maketitle

\noindent\textbf{Status of this note.}
\begin{itemize}
\item \emph{Authorship.}  Written on 23 November 2026 by an AI agent (Claude, the drafter
``route-t'' of the project's working council).  It was commissioned by the author and supervised
by the council's orchestrator (an AI agent), as the content proposed for Phase C3 of
\path{ugjs/PLAN.md}.
\item \emph{Audits.}  On 23 November 2026 two independent agents audited it, and each re-checked
the revision; their reports are \path{ugjs/council/route-t-audit-analysis.md} and
\path{ugjs/council/route-t-audit-sources.md}.  No finding was CRITICAL or MAJOR.  The note was
revised for the findings of the audits, and a last round of MINOR points from the re-checks was
applied after them.  One of the two auditors read that last round in full on 24 November 2026 and
confirmed it.  Of the three points of that confirmation (two MINOR, one cosmetic; section 10 of
the first report), the second, on the record of 23 November, is taken up in the item ``Which
record'' below; the other two are not applied: Simon's class $F$ on p.~163 is fixed by his own
notation and is not a reading, as the proof of Lemma~T.15 calls it; and three wordings.  The
sentences of this item and of ``Which record'' that concern the confirmation were written by the
orchestrating agent and read only by the documentation check of 24 November
(\path{ugjs/council/route-t-docs-check.md}).
\item \emph{Not read by a human.}  Nothing in it has yet been read by a human.
\item \emph{Not integrated.}  This note changes nothing in the record: it is not integrated, and
Theorem~12.2 of \path{part_i/main.tex} is stated and proved there as before.
\item \emph{Which record.}  The record it refers to is \path{part_i/main.tex} at commit 49f2d23;
line numbers and pages are those of that commit.  The record of 23 November 2026 (commit
9733263, 127 pages) has the same statement numbers as 49f2d23, and different pages; it rewords one
sentence of the imports paragraph of Theorem~12.2.  It also restricts Corollary~12.3 to
$\mathscr P(\xi)=\xi^4$ and to cutoffs $0\le g\le1$: its audit found that the statement at
49f2d23, which covered every cutoff of GJ3, goes beyond its proof, which needs $0\le g\le1$.
This note obtains Corollary~12.3 from that proof with $\Psit_{m_1}$ in place of $\Psi_{m_1}$
(\S\ref{sec:assembly}(b)), so it gives the restricted corollary only, and Corollary~12.3 is to
be read so wherever this note names it.  By the auditor's confirmation (point C2), the argument
applies to the restricted corollary, whose proof uses Theorem~12.2(1) only through its rate.  The
restricted corollary also rests on the reading of A1's Scope applied to GJ3, which the record of
23 November names after it.
\end{itemize}

\section*{Summary}

\emph{What is proved.}  In the record (\path{part_i/main.tex} at commit 49f2d23), items
(1)--(4) of Theorem~12.2 (label \texttt{thm:MAIN}) rest on the uniform local moment $M$ of
Theorem~5.1 (\texttt{thm:M2}), which rests on Theorem~D.4 (\texttt{thm:D4}) and through it on the
transverse case A8(iv)-T of GJS Theorem~1.1.7.  This note replaces, in the chain from the
uniform gap to items (1)--(4), the moment $M$ by the resolvent-weighted moment
$\Mt(\lambda_0)=\sup\|(\lambda_0-iH(g))^{-1}Q_gV(h)\Om_g\|$ (Definition~\ref{def:Mt}), proves
the modified chain (Lemmas~\ref{lem:transforms}--\ref{lem:rate}), and bounds $\Mt(\lambda_0)$
for every $\lambda_0\ge1$ by $c_k^{1/2}M_B$, $c_k=2/(1+e^{-2})$, through a time-integrated copy
of Theorem~D.4 (Theorems~\ref{thm:TD4} and~\ref{thm:Mt}).  Items (1)--(4) of Theorem~12.2 and
Corollary~12.3 then hold with the same state $\omega_\infty$ and a new explicit rate function
$\Psit_{m_1}$, decaying faster than every inverse power (Theorem~\ref{thm:routeT}).  These proofs
use neither Theorem~5.1 nor Theorem~D.4 nor A8(iv)-T.  The sketch of
22 November (steps (a)--(d)) is confirmed; its constants are confirmed or slightly improved.

\emph{What it rests on, honestly.}  Route T adds no input to the record's imports.  Its
Euclidean-to-Hamiltonian step (Lemma~\ref{lem:fkn}) uses A2's multi-time formula for bounded
functions at two times placed asymmetrically in a slab; formulas (2) and (3) of the record's
Lemma~4.1 place two symmetric times and one time.  The note derives that formula from the
bounded case of Simon's Theorem~V.15 as printed, by approximation, and passes to the unbounded
$V(h)$ by the truncation of the record's Lemma~4.3(ii).  Every
application of GJS Theorem~1.1.7 that the new chain makes pairs two time-integrated monomials:
a two-factor piece of $\Vb(h')^2$ and a one-square interaction $U_{\Delta,n}$, across strips of
both orientations (cases P3, P4 of \S\ref{sec:cases}).  The printed statement covers these
pairs.  The printed proof reaches them on reading R1 of p.~596 (\S\ref{sec:readings}): the
monomial $A$ of Theorem~1.1.11 may have its factors at different times.  For the transverse
case P4 this holds granted the rotation covariance that the proof's ``rotating'' presupposes.
R1 is the reading under which GJS's own proofs are complete: their application of
Theorem~1.1.7 on p.~595, and their derivation of Theorem~1.1.8 at its printed multi-time scope
from Theorem~1.1.11 (pp.~600, 629).  No printed sentence states R1 in words.  On the other
reading (R2) the new applications are not reached, the parallel ones included (only the trivial
pieces with a factor of degree $0$ are).  Remark~\ref{rem:obstruction} argues that no variant
that bounds a quadratic observable by the comparison of Lemma~D.2 around a box of fixed shape
can be reached on R2.  So Route T does not make items (1)--(4) rest only on cases the printed
proof reaches.  It trades A8(iv)-T, which no reading of the printed proof reaches, for
reading R1.  The moment clause of item (5), Theorem~5.1 and Corollary~5.2 are not touched and
still rest on A8(iv)-T.

\emph{Remark N.39.}  The finding N3 of 22 November is confirmed: the parallel pairs of
Theorem~D.3 are reached after slicing $U_{\Delta,n}$ in time, by a Fubini step written out in
Proposition~\ref{prop:N39}, on either reading of p.~596, with no multi-time reading.

\tableofcontents

\section{Setting and what is to be proved}\label{sec:setting}

\subsection{Sources and conventions}

\emph{The record} is \path{part_i/main.tex} at commit 49f2d23 (8744 lines).  Statements of the
record are cited by their number, label and page as read from \path{part_i/main.aux}, for
example ``Lemma~9.3 (\texttt{lem:cluster}, p.~37)''; line numbers are those of that commit.
The record's Appendix~D numbers its auxiliary statements N.1--N.55; they are cited the same
way.  Assumptions A1--A8 are the record's standing imports (\S2.4).  \emph{GJS} is
J.~Glimm, A.~Jaffe, T.~Spencer, \emph{The Wightman axioms and particle structure in the
$\mathscr P(\phi)_2$ quantum field model}, Annals of Mathematics 100 (1974), no.~3, 585--632
(checked on the first page of the copy in \path{refs/}, which is not part of the repository);
page images of pp.~593--597 are \path{audit_2026/evidence/gjs-theorems-10.png} (p.~593),
\path{gjs-theorems-11.png} (p.~594), \path{gjs-theorems-12.png} (p.~595),
\path{gjs-thm1111-13.png} (p.~596) and \path{gjs-proof-14.png} (p.~597), PDF page indices
10--14 of the copy in \path{refs/}; pp.~600, 628 and 629 are \path{gjs-proof118-17.png},
\path{gjs-thm41-45.png} and \path{gjs-uniform-46.png} (PDF page indices 17, 45, 46), and
p.~590 is PDF page index 7.  \emph{Simon} is B.~Simon, \emph{The $P(\Phi)_2$
Euclidean (Quantum) Field Theory}, Princeton University Press, 1974; its Theorem~V.15 is on
printed p.~163 (image \path{audit_2026/evidence/simon-fkn-185.png}) and continues on p.~164
(PDF page index 186 of the copy in \path{refs/}).  Every quotation of a source below was read
on the rendered page on 23 November 2026.  \emph{The commission} below is the author's
commission of 23 November 2026 for this note (Phase C3 of \path{ugjs/PLAN.md}).  It asked for
items (a)--(d), which follow steps (a)--(d) of the sketch of 22 November, and for the three
categories of \S\ref{sec:input}.

\emph{Standing assumption.}  As in \S\S8--10 of the record, $0\le\lambda/m_0^2<\varepsilon_{\mathscr
P}$, and we work in the working coordinates of \S2.6 (\texttt{subsec:weak}, p.~22) and
Remark~3.3 (\texttt{rem:wlog}, p.~24), in which the printed hypotheses of GJS Theorem~1.1.7
hold.  $\gamma>0$ is a constant for which (G$_\gamma$) of the record (\texttt{eq:Ggamma},
p.~34) holds:
\[
 \spec H(g)\subseteq\{0\}\cup[\gamma,\infty),\qquad\dim\ker H(g)=1\qquad(g\in\CH);
\]
Theorem~7.8 (\texttt{thm:G}, p.~33) gives $\gamma=m_1$.  For $g\in\CH$ we write
$H=H(g)$, $\Om=\Om_g$, $P_0=|\Om\rangle\langle\Om|$, $Q=1-P_0$, and $P_H(\cdot)$ for the
spectral measure of $H$; under (G$_\gamma$), $Q=P_H([\gamma,\infty))$ and $H^{-1}Q$ denotes
the bounded operator $\phi(H)$, $\phi(E)=E^{-1}\mathbf 1_{[\gamma,\infty)}(E)$, of norm
$\le\gamma^{-1}$ (as in Proposition~8.2, \texttt{prop:projreg}, p.~35).  The path setting of
\S8 of the record is used verbatim: $g,g'\in\G$ equal to $1$ on a common nonempty open
interval, $g_s=(1-s)g+sg'$, $v=g'-g$, $H_s$, $\Om_s$, $P_s$, $Q_s$; the partition of unity
$\eta_j$, $h_j=\eta_jv$, $V_j=V(h_j)$, $B_j=\supp h_j$ of Lemma~10.2 (\texttt{lem:localize},
p.~38); for a bounded open interval $O$, $r_j=\dist(B_j,O)$.  $W_\gamma$ and
$T_\gamma(r)=\int_{|t|>r}|W_\gamma|$ are those of Lemma~9.1 (\texttt{lem:filter}, p.~35).

The resolvent parameter is a number $\lambda_0\ge1$; it has nothing to do with the coupling
constant $\lambda$.

\begin{tdef}[the resolvent-weighted moment]\label{def:Mt}
Let $\mathscr A$ be the set of pairs $(g,h)$ with $g\in\CH$ and $h$ real and measurable,
$\|h\|_\infty\le1$, $\supp h\subseteq[a,a+2]$ for some $a\in\R$.  For $\lambda_0\ge1$ let
\[
 \Mt(\lambda_0):=\sup\Bigl\{\bigl\|(\lambda_0-iH(g))^{-1}Q_g\,V(h)\,\Om_g\bigr\|:\
 (g,h)\in\mathscr A\Bigr\}\in[0,\infty] .
\]
Here $Q_g=1-|\Om_g\rangle\langle\Om_g|$, and $\Om_g\in D(V(h))$ by Proposition~2.6
(\texttt{prop:smoothing}, p.~22).  The operator $(\lambda_0-iH(g))^{-1}$ is bounded with norm
$\le\lambda_0^{-1}$ because $H(g)$ is self-adjoint and $|\lambda_0-iE|\ge\lambda_0$ for real
$E$.  The record's moment of Theorem~5.1 is the supremum of
$\|V(h)\Om_g\|$ over the same set; since $|\lambda_0-iE|\ge\lambda_0\ge1$,
$\Mt(\lambda_0)\le M$ whenever the record's bound $M$ is available.  Route T does not use it.
\end{tdef}

\subsection{The modified chain}\label{sec:chain}

The following table says which statements of the record are replaced in the chain
``uniform gap $\Rightarrow$ items (1)--(4)''.  Statements not listed are used as they stand.

\begin{center}\small
\begin{tabular}{@{}p{0.36\textwidth}p{0.10\textwidth}p{0.46\textwidth}@{}}
\toprule
record (label, page) & uses $M$? & in Route T\\
\midrule
Lemma 9.2 (\texttt{lem:separation}, p.~36) & no & kept; Lemma~\ref{lem:sep} is its variant
with an exponential discrepancy\\
Lemma 9.3 (\texttt{lem:cluster}, p.~37) & no & kept (it is used in \S13); Lemma~\ref{lem:cluster}
is the resolvent-weighted variant\\
Lemma 10.1 (\texttt{lem:statederiv}, p.~38) & no & replaced in the chain by
Lemma~\ref{lem:deriv}\\
Lemma 10.2 (\texttt{lem:localize}, p.~38) & yes & only its first assertion (finite
decomposition) is used, in Lemma~\ref{lem:deriv}\\
Lemma 10.3 (\texttt{lem:tworegime}, p.~38) & yes & replaced by Lemma~\ref{lem:tworegime}\\
Proposition 10.4 (\texttt{prop:onepath}, p.~38) & yes & replaced by
Proposition~\ref{prop:onepath}\\
Lemma 10.5 (\texttt{lem:rate}, p.~39) & yes & replaced by Lemma~\ref{lem:rate}\\
Theorem 5.1 (\texttt{thm:M2}, p.~27) & --- & not used; Theorem~\ref{thm:Mt} bounds
$\Mt(\lambda_0)$\\
Theorem D.4 (\texttt{thm:D4}, p.~100) & --- & not used; Theorem~\ref{thm:TD4} is its
time-integrated copy\\
Theorem 11.1, Proposition 11.2 (p.~40), Theorem 12.2 (1)--(4), Corollary 12.3 (pp.~41--42) &
through $\Psi_{m_1}$ & same statements with $\Psit_{m_1}$ (Theorem~\ref{thm:routeT})\\
Theorem 12.2 (5) & --- & unchanged; its moment clause still rests on Theorem~5.1, hence on
A8(iv)-T\\
\bottomrule
\end{tabular}
\end{center}

The constants of the new chain are $c_k:=2/(1+e^{-2})$ ($=1.7615\ldots$) and
\begin{gather}
 k(E):=\iint_{[-\frac12,\frac12]^2}e^{-|\tau-\tau'|E}\,d\tau\,d\tau'
 =\frac{2(E-1+e^{-E})}{E^2}\quad(E>0),\qquad k(0)=1,\label{eq:k}\\
 \Xi_{\gamma,\lambda_0}(R):=\Bigl(1+\frac{4}{\pi\gamma R}\Bigr)e^{-\gamma R/4}
 +\Bigl(1+\frac1\pi\Bigr)e^{-\lambda_0R/2}\qquad(R>0),\label{eq:Xi}\\
 \begin{split}
 \Psit_{\gamma,\lambda_0}(d):=12\,M'\sum_{n\ge\lfloor d\rfloor}\Bigl(\Xi_{\gamma,\lambda_0}(n)
 &+\lambda_0\|W_\gamma\|_1\,\Xi_{\gamma,\lambda_0}(n/2)\\
 &+\lambda_0\,T_\gamma(n/2)\Bigr)\qquad(d\ge1),
 \end{split}\label{eq:Psit}
\end{gather}
and $\Psit_{\gamma,\lambda_0}(d):=\max\{2,\Psit_{\gamma,\lambda_0}(1)\}$ for $0\le d<1$, where
$M'$ is any number with $M'\ge\Mt(\lambda_0)$.  In Theorem~\ref{thm:routeT},
$M'=\Mt_*:=c_k^{1/2}M_B$, with $M_B$ the constant of Theorem~D.4 (equal to the $M$ of
Theorem~5.1), and $\gamma=m_1$.

\begin{ttheorem}[Route T]\label{thm:routeT}
Let $0\le\lambda/m_0^2<\varepsilon_{\mathscr P}$ and $\lambda_0\ge1$.  Put
$\Psit_{m_1}:=\Psit_{m_1,\lambda_0}$ with $M'=\Mt_*=c_k^{1/2}M_B$.
\begin{enumerate}
\item[\rm(a)] $\Mt(\lambda_0)\le\Mt_*$.
\item[\rm(b)] Theorem~11.1 (\texttt{thm:cauchy}) and Proposition~11.2 (\texttt{prop:normal}) of
the record hold with $\Psi_{m_1}$ replaced by $\Psit_{m_1}$.  Consequently items (1)--(4) of
Theorem~12.2 and Corollary~12.3 hold, for the same state $\omega_\infty$, with $\Psi_{m_1}$
replaced by $\Psit_{m_1}$ in item (1) and in the proof of Corollary~12.3; $\Psit_{m_1}$ decays
faster than every inverse power (Lemma~\ref{lem:rate}).
\item[\rm(c)] The proofs of (a) and (b) use neither Theorem~5.1 nor Theorem~D.4 nor the item
A8(iv)-T of Assumption~A8.  Every application of GJS Theorem~1.1.7 they make is one of the
cases P1--P4 of \S\ref{sec:cases}: P1 and P2 are those of Lemma~7.4 and Theorem~D.3 (strips
parallel to the time-zero line), P3 and P4 are those of Theorem~\ref{thm:TD4} (two
time-integrated monomials, parallel and transverse strips).
\end{enumerate}
\end{ttheorem}

\emph{Constants.}  $\Mt_*$ depends on $(\lambda,\mathscr P,m_0)$ through $M_B$, i.e.\ through
$C_7$, $m_1$, $K_1$, $K_2$, $V_*$ and the unevaluated $S_2(B_{3\times5})$; it does not depend on
$\lambda_0\ge1$, $g$ or $h$.  $\Psit_{m_1}$ depends in addition on $\lambda_0$ and on the
cut-off function $\chi$ of Lemma~9.1 (through $W_{m_1}$ and $T_{m_1}$), and on nothing else; it
is uniform in $g,g'\in\G$, in $O$ and in $A$.  The threshold of the one-path estimate is
$d\ge1$, in place of $d\ge d_0=4+2/\gamma$.  Any $\lambda_0\ge1$ may be used; $\lambda_0=1$ is
the simplest choice, and $\lambda_0=\max\{1,m_1/2\}$ makes the two exponentials of
$\Xi_{m_1,\lambda_0}$ decay at the same rate $e^{-m_1R/4}$ when $m_1\ge2$.

\emph{Coordinates.}  Theorem~\ref{thm:routeT}, like Theorem~12.2, is proved in the working
coordinates.  Remark~3.3(1) of the record (ll.~1660--1669 and 1698--1701) transports item (1)
to a general pair $(m_0,\lambda)$ through the dilation unitary and uses of the
working-coordinate statement only its inequality; the same transport applies verbatim to
$\Psit$.  In the original coordinates the rate is $d\mapsto\Psit_{m_{1,*},\lambda_0}(s\,d)$,
$s=m_0/M_*$, with $m_{1,*}$ the gap of the working coordinates and $\lambda_0$ in working units.

\emph{Status.}  Proved, conditional on: the imports A1--A8 of the record, with A8(iv) used only
in the cases P1--P4 and A8(iv)-T not used, and with A2 used at two times placed asymmetrically
in a slab (Lemma~\ref{lem:fkn}(ii)), as formula (3) of the record's Lemma~4.1 uses it at one; the readings
the record declares (Scope paragraphs of A1, A3, A6, and A8 Scope (d), with the open items of
\path{part_i/C5_reading_record.md}); and the statements of the record listed in
\S\ref{sec:status}.  Whether the printed \emph{proof} of Theorem~1.1.7 reaches P3--P4 depends
on reading R1 and, for P4, on the rotation covariance the proof presupposes
(\S\ref{sec:input}).  The proof of Theorem~\ref{thm:routeT} is in \S\ref{sec:assembly}.

\section{The Hamiltonian side: steps (a) and (b) of the sketch}\label{sec:hamiltonian}

\begin{tlemma}[resolvent transforms]\label{lem:transforms}
Let $H\ge0$ be self-adjoint with $\spec H\subseteq\{0\}\cup[\gamma,\infty)$, $\gamma>0$, and
$\ker H$ spanned by a unit vector $\Om$; put $P_0=|\Om\rangle\langle\Om|$, $Q=1-P_0$.  Let
$X=X^*$ with $\Om\in D(X)$, let $A$ be bounded, $\lambda_0>0$, and put
\[
 u:=(\lambda_0-iH)^{-1}QX\Om,\qquad \mt:=\|u\|\,\|A\|,
\]
\begin{gather*}
 F(t):=\ip{X\Om}{(e^{itH}-P_0)A\Om},\qquad G(t):=\ip{A^*\Om}{(e^{-itH}-P_0)X\Om},\\
 \Ft(t):=\int_{-\infty}^te^{-\lambda_0(t-s)}F(s)\,ds,\qquad
 \Gt(t):=\int_{-\infty}^te^{-\lambda_0(t-s)}G(s)\,ds .
\end{gather*}
\begin{enumerate}
\item[\rm(i)] $F$ and $G$ are continuous and bounded by $\|X\Om\|\,\|A\|$; the integrals
defining $\Ft$, $\Gt$ converge absolutely.
\item[\rm(ii)] $\Ft(t)=\ip{u}{e^{itH}QA\Om}=\int e^{itE}\,d\tilde\rho(E)$ with the complex
Borel measure $\tilde\rho(\cdot):=\ip{u}{P_H(\cdot)QA\Om}$, which is supported in
$[\gamma,\infty)$ and has total variation $|\tilde\rho|(\R)\le\mt$.
\item[\rm(iii)] $\Gt(t)=\ip{QA^*\Om}{e^{-itH}u}=\int e^{itE}\,d\tilde\sigma(E)$, where
$\tilde\sigma$ is the image of $\ip{QA^*\Om}{P_H(\cdot)u}$ under $E\mapsto-E$; it is supported
in $(-\infty,-\gamma]$ and $|\tilde\sigma|(\R)\le\mt$.
\item[\rm(iv)] $|\Ft(t)|\le\mt$ and $|\Gt(t)|\le\mt$ for every $t\in\R$.
\end{enumerate}
\end{tlemma}

\begin{proof}
(i) $HP_0=0$ gives $e^{itH}P_0=P_0$, so $e^{itH}-P_0=e^{itH}Q$ and
$F(t)=\ip{X\Om}{e^{itH}QA\Om}$, $G(t)=\ip{A^*\Om}{e^{-itH}QX\Om}$.  Strong continuity of
$t\mapsto e^{itH}$ gives continuity, and $\|QA\Om\|\le\|A\|$ gives the bound; $e^{-\lambda_0(t-s)}$
is integrable over $s\in(-\infty,t]$.

(ii) For every vector $y$, $v\mapsto e^{-\lambda_0v}e^{-ivH}y$ is continuous and norm-integrable
on $[0,\infty)$, and for every vector $x$, by Fubini (the scalar measure
$\ip{x}{P_H(\cdot)y}$ is finite and $e^{-\lambda_0v}$ is integrable),
\begin{align*}
 \Bigl\langle x,\int_0^\infty e^{-\lambda_0v}e^{-ivH}y\,dv\Bigr\rangle
 &=\int\Bigl(\int_0^\infty e^{-\lambda_0v}e^{-ivE}dv\Bigr)d\ip{x}{P_H(E)y}\\
 &=\int\frac{d\ip{x}{P_H(E)y}}{\lambda_0+iE}=\ip{x}{(\lambda_0+iH)^{-1}y} .
\end{align*}
Hence, substituting $s=t-v$,
\[
 \Ft(t)=\int_0^\infty e^{-\lambda_0v}\ip{X\Om}{e^{i(t-v)H}QA\Om}\,dv
 =\ip{X\Om}{(\lambda_0+iH)^{-1}e^{itH}QA\Om} .
\]
The vector $e^{itH}QA\Om$ lies in $Q\HH$, and $(\lambda_0+iH)^{-1}$ commutes with $Q$, so
$(\lambda_0+iH)^{-1}e^{itH}QA\Om=Q(\lambda_0+iH)^{-1}e^{itH}QA\Om$; with
$((\lambda_0+iH)^{-1})^*=(\lambda_0-iH)^{-1}$ this gives
$\Ft(t)=\ip{(\lambda_0-iH)^{-1}QX\Om}{e^{itH}QA\Om}=\ip{u}{e^{itH}QA\Om}$, and the spectral
theorem gives $\Ft(t)=\int e^{itE}d\tilde\rho(E)$.  Under the hypothesis on $\spec H$,
$Q=P_H([\gamma,\infty))$, so $\tilde\rho(B)=\ip{u}{P_H(B\cap[\gamma,\infty))QA\Om}$ and
$\supp\tilde\rho\subseteq[\gamma,\infty)$.  For a finite Borel partition $(B_i)$ of $\R$ and
$w:=QA\Om$, Cauchy--Schwarz twice gives
\begin{align*}
 \sum_i|\tilde\rho(B_i)|=\sum_i|\ip{P_H(B_i)u}{P_H(B_i)w}|
 &\le\Bigl(\sum_i\|P_H(B_i)u\|^2\Bigr)^{1/2}\Bigl(\sum_i\|P_H(B_i)w\|^2\Bigr)^{1/2}\\
 &=\|u\|\,\|w\|\le\mt ;
\end{align*}
the supremum over partitions is $|\tilde\rho|(\R)$.

(iii) In the same way, with $\int_0^\infty e^{-\lambda_0v}e^{ivE}dv=(\lambda_0-iE)^{-1}$,
\begin{align*}
 \Gt(t)&=\int_0^\infty e^{-\lambda_0v}\ip{A^*\Om}{e^{-i(t-v)H}QX\Om}\,dv\\
 &=\ip{A^*\Om}{e^{-itH}(\lambda_0-iH)^{-1}QX\Om}=\ip{QA^*\Om}{e^{-itH}u},
\end{align*}
since $e^{-itH}u\in Q\HH$.  The measure $\sigma'(\cdot):=\ip{QA^*\Om}{P_H(\cdot)u}$ is supported
in $[\gamma,\infty)$ and has variation $\le\|QA^*\Om\|\,\|u\|\le\mt$ (the computation of (ii));
$\Gt(t)=\int e^{-itE}d\sigma'(E)=\int e^{itE}d\tilde\sigma(E)$ for its image $\tilde\sigma$
under $E\mapsto-E$.

(iv) $|\Ft(t)|\le|\tilde\rho|(\R)$ and $|\Gt(t)|\le|\tilde\sigma|(\R)$.
\end{proof}

\begin{tlemma}[the locality discrepancy]\label{lem:discrepancy}
In the setting of Lemma~\ref{lem:transforms}, suppose $F(t)=G(t)$ for all $t\in(-R,R)$, where
$R>0$.  Then, for every $t\in(-R,R)$,
\[
 \Ft(t)-\Gt(t)=c\,e^{-\lambda_0t},\qquad c:=\int_{-\infty}^{-R}e^{\lambda_0s}\bigl(F(s)-G(s)\bigr)ds,
 \qquad |c|\le2\,\mt\,e^{-\lambda_0R} .
\]
\end{tlemma}

\begin{proof}
For $t\in(-R,R)$, $F-G$ vanishes on $(-R,t]$, so
$\Ft(t)-\Gt(t)=\int_{-\infty}^{-R}e^{-\lambda_0(t-s)}(F-G)(s)\,ds=e^{-\lambda_0t}c$; the integral
defining $c$ converges absolutely by Lemma~\ref{lem:transforms}(i).  By
Lemma~\ref{lem:transforms}(iv), $|c|e^{-\lambda_0t}\le2\mt$, i.e.\ $|c|\le2\mt e^{\lambda_0t}$,
for every $t\in(-R,R)$; letting $t\downarrow-R$ gives the bound.
\end{proof}

So the resolvent smoothing destroys the exact equality $F=G$ on $(-R,R)$ that Lemma~9.2 uses,
but leaves a discrepancy that is explicit, smooth and exponentially small on $(-R/2,R/2)$:
there $|ce^{-\lambda_0t}|\le2\mt e^{-\lambda_0R/2}$.  The next lemma is Lemma~9.2 with this
discrepancy; the principal-value term it produces is computed exactly.

\begin{tlemma}[spectral separation with an exponential discrepancy]\label{lem:sep}
Let $\gamma>0$, $\lambda_0>0$, $R>0$, $0<R'\le R$ and $\mt\ge0$.  Let $\tilde\rho$, $\tilde\sigma$
be complex Borel measures on $\R$ with $\supp\tilde\rho\subseteq[\gamma,\infty)$,
$\supp\tilde\sigma\subseteq(-\infty,-\gamma]$ and $|\tilde\rho|(\R),|\tilde\sigma|(\R)\le\mt$;
put $\Ft(t)=\int e^{itE}d\tilde\rho(E)$, $\Gt(t)=\int e^{itE}d\tilde\sigma(E)$, and suppose
$\Ft(t)-\Gt(t)=c\,e^{-\lambda_0t}$ for all $t\in(-R,R)$, with $c\in\C$.  Then
\begin{equation}
 |\Ft(0)|\le\mt\Bigl(1+\frac{2}{\pi\gamma R'}\Bigr)e^{-\gamma R'/2}
 +|c|\Bigl(\frac12+\frac{\sinh(\lambda_0R')}{\pi}\Bigr).
 \label{eq:sep}
\end{equation}
\end{tlemma}

\begin{proof}
\emph{Step 1: the separator} (Lemma~9.2, Step~1, re-derived).  For $a,\varepsilon>0$ let
$f_{a,\varepsilon}(t):=(2\pi i)^{-1}e^{-t^2/(2a)}(t-i\varepsilon)^{-1}$; it is integrable, since
$|t-i\varepsilon|\ge\varepsilon$.  From $i\int_0^\infty e^{-(\varepsilon+it)E'}dE'=
i(\varepsilon+it)^{-1}=(t-i\varepsilon)^{-1}$,
\[
 f_{a,\varepsilon}(t)=\frac1{2\pi}\int_0^\infty e^{-\varepsilon E'}e^{-t^2/(2a)-itE'}\,dE' ,
\]
and Fubini (the integrand of the double integral is absolutely integrable) together with
$\int e^{-t^2/(2a)}e^{itk}dt=(2\pi a)^{1/2}e^{-ak^2/2}$ gives
\[
 \check f_{a,\varepsilon}(E):=\int f_{a,\varepsilon}(t)e^{itE}dt
 =\int_0^\infty\Bigl(\frac a{2\pi}\Bigr)^{1/2}e^{-a(E-E')^2/2}e^{-\varepsilon E'}\,dE'\in[0,1],
\]
which increases, as $\varepsilon\downarrow0$, to $\Phi(\sqrt aE)$, $\Phi$ the standard normal
distribution function (monotone convergence).  This is the transform (\texttt{eq:septransform},
(20)) of the record.

\emph{Step 2: pairing and tail.}  By Fubini ($f_{a,\varepsilon}\in L^1$, finite measures),
$\int f_{a,\varepsilon}\Ft\,dt=\int\check f_{a,\varepsilon}d\tilde\rho$ and
$\int f_{a,\varepsilon}\Gt\,dt=\int\check f_{a,\varepsilon}d\tilde\sigma$.  For $|t|\ge R'$,
$|f_{a,\varepsilon}(t)|\le(2\pi)^{-1}e^{-t^2/(2a)}|t|^{-1}\le(2\pi)^{-1}e^{-t^2/(2a)}|t|/R'^2$,
so, uniformly in $\varepsilon$,
\[
 \int_{|t|\ge R'}|f_{a,\varepsilon}(t)|\,dt\le\frac1{\pi R'^2}\int_{R'}^\infty te^{-t^2/(2a)}dt
 =\frac{a}{\pi R'^2}e^{-R'^2/(2a)} .
\]

\emph{Step 3: the principal-value term.}  Let
\[
 \varphi(t):=e^{-t^2/(2a)}e^{-\lambda_0t},\qquad
 J:=\int_0^{R'}e^{-t^2/(2a)}\,\frac{\sinh(\lambda_0t)}{t}\,dt .
\]
We show
\begin{equation}
 \lim_{\varepsilon\downarrow0}\int_{-R'}^{R'}f_{a,\varepsilon}(t)\,c\,e^{-\lambda_0t}\,dt
 =\frac c2+\frac{ic}{\pi}J,\qquad
 \Bigl|\frac c2+\frac{ic}\pi J\Bigr|\le|c|\Bigl(\frac12+\frac{\sinh(\lambda_0R')}{\pi}\Bigr).
 \label{eq:pv}
\end{equation}
Write $(t-i\varepsilon)^{-1}=t(t^2+\varepsilon^2)^{-1}+i\varepsilon(t^2+\varepsilon^2)^{-1}$.
With $t=\varepsilon s$, $\varepsilon\int_{-R'}^{R'}\varphi(t)(t^2+\varepsilon^2)^{-1}dt=
\int_{-R'/\varepsilon}^{R'/\varepsilon}\varphi(\varepsilon s)(1+s^2)^{-1}ds\to\pi\varphi(0)=\pi$
by dominated convergence (dominator $\sup_{[-R',R']}|\varphi|\,(1+s^2)^{-1}$).  Next,
$\int_{-R'}^{R'}\varphi(t)t(t^2+\varepsilon^2)^{-1}dt=\int_0^{R'}(\varphi(t)-\varphi(-t))
t(t^2+\varepsilon^2)^{-1}dt\to\int_0^{R'}(\varphi(t)-\varphi(-t))t^{-1}dt$ by dominated
convergence, the integrand being bounded by $|\varphi(t)-\varphi(-t)|/t\le2\sup_{[-R',R']}
|\varphi'|$ (mean value theorem).  Since $\varphi(t)-\varphi(-t)=-2e^{-t^2/(2a)}\sinh(\lambda_0t)$,
the limit of $\int_{-R'}^{R'}\varphi(t)(t-i\varepsilon)^{-1}dt$ is $-2J+i\pi$, and multiplying by
$c/(2\pi i)$ gives $c/2-cJ/(\pi i)=c/2+icJ/\pi$.  As $J$ is real,
$|c/2+icJ/\pi|=|c|(\frac14+J^2/\pi^2)^{1/2}\le|c|(\frac12+J/\pi)$.  Finally
$e^{-t^2/(2a)}\le1$ and $\sinh(x)/x\le\cosh(x)$ for $x>0$ (that is, $\tanh x\le x$) give
$J\le\int_0^{\lambda_0R'}\sinh(x)x^{-1}dx\le\int_0^{\lambda_0R'}\cosh(x)\,dx=\sinh(\lambda_0R')$.

\emph{Step 4: combination.}  Since $R'\le R$, the hypothesis gives
$\Ft-\Gt=ce^{-\lambda_0t}$ on $(-R',R')$, so for every $\varepsilon>0$
\[
 \int\check f_{a,\varepsilon}d\tilde\rho-\int\check f_{a,\varepsilon}d\tilde\sigma
 =\int_{|t|\ge R'}f_{a,\varepsilon}(\Ft-\Gt)\,dt+\int_{-R'}^{R'}f_{a,\varepsilon}(t)\,ce^{-\lambda_0t}dt .
\]
As $\varepsilon\downarrow0$ the left side tends to $\int\Phi(\sqrt aE)\,d\tilde\rho-\int\Phi(\sqrt
aE)\,d\tilde\sigma$ (dominated convergence, $0\le\check f_{a,\varepsilon}\le1$), and the last
term tends to the limit (\ref{eq:pv}); hence the first term on the right converges too, and
since $|\Ft-\Gt|\le2\mt$ everywhere, Step~2 bounds it by $2\mt a(\pi R'^2)^{-1}e^{-R'^2/(2a)}$.
For $x\ge0$, $1-\Phi(x)=\int_0^\infty(2\pi)^{-1/2}e^{-(x+y)^2/2}dy\le e^{-x^2/2}/2$.  On
$\supp\tilde\rho$, $E\ge\gamma$ and $0\le1-\Phi(\sqrt aE)\le e^{-a\gamma^2/2}/2$; on
$\supp\tilde\sigma$, $E\le-\gamma$ and $0\le\Phi(\sqrt aE)\le e^{-a\gamma^2/2}/2$.  Writing
\[
 \Ft(0)=\int d\tilde\rho=\int\bigl(1-\Phi(\sqrt aE)\bigr)d\tilde\rho
 +\Bigl[\int\Phi(\sqrt aE)\,d\tilde\rho-\int\Phi(\sqrt aE)\,d\tilde\sigma\Bigr]
 +\int\Phi(\sqrt aE)\,d\tilde\sigma
\]
we obtain $|\Ft(0)|\le\mt e^{-a\gamma^2/2}+2\mt a(\pi R'^2)^{-1}e^{-R'^2/(2a)}
+|c|(\frac12+\sinh(\lambda_0R')/\pi)$.  The choice $a:=R'/\gamma$ makes
$a\gamma^2/2=R'^2/(2a)=\gamma R'/2$ and $a/(\pi R'^2)=1/(\pi\gamma R')$, which is
(\ref{eq:sep}).
\end{proof}

\emph{Boundary values.}  For $c=0$ and $R'=R$, (\ref{eq:sep}) is the bound of Step~4 of the
record's Lemma~9.2 before its simplification by $\gamma R\ge1$.  As $R'\downarrow0$ the bound
becomes useless; the trivial bound $|\Ft(0)|\le\mt$ always holds.  The test of Step~3 on
three parameter sets is in Appendix~\ref{app:numerics} (\emph{numerical}).

\begin{tlemma}[resolvent-weighted clustering; replaces Lemma 9.3 in the chain]\label{lem:cluster}
Let $g\in\CH$, $H=H(g)$, $\Om=\Om_g$, and assume (G$_\gamma$).  Let $B\subset\R$ be compact, let
$X=X^*$ be affiliated with $\A(O')'$ for every bounded open interval $O'$ with $\dist(B,O')>0$,
with $\Om\in D(X)$, and let $A\in\A(O)$ for a bounded interval $O$ with $R:=\dist(B,O)>0$.
Let $\lambda_0>0$ and $\mt:=\|(\lambda_0-iH)^{-1}QX\Om\|\,\|A\|$.  Then
\[
 \bigl|\ip{(\lambda_0-iH)^{-1}QX\Om}{QA\Om}\bigr|\le\mt\,\min\{1,\Xi_{\gamma,\lambda_0}(R)\} .
\]
\end{tlemma}

\begin{proof}
The functions $F$, $G$ of Lemma~\ref{lem:transforms} are the functions $F$, $G$ of the proof of
Lemma~9.3.  Step~2 of that proof shows $F=G$ on $(-R,R)$; it uses Proposition~2.7
(\texttt{prop:propagation}), the affiliation of $X$, $\Om\in D(X)$, $X=X^*$ and
$e^{-itH}\Om=\Om$, and not the hypothesis $\gamma R\ge1$ of Lemma~9.3.  By
Lemma~\ref{lem:transforms}(ii), $\Ft(0)=\ip{(\lambda_0-iH)^{-1}QX\Om}{QA\Om}$, and the
measures there satisfy the hypotheses of Lemma~\ref{lem:sep}; Lemma~\ref{lem:discrepancy} gives
$c$ with $|c|\le2\mt e^{-\lambda_0R}$.  Lemma~\ref{lem:sep} with $R'=R/2$ gives
\begin{align*}
 |\Ft(0)|&\le\mt\Bigl(1+\frac4{\pi\gamma R}\Bigr)e^{-\gamma R/4}
 +2\mt e^{-\lambda_0R}\Bigl(\frac12+\frac{\sinh(\lambda_0R/2)}\pi\Bigr)\\
 &\le\mt\Bigl(1+\frac4{\pi\gamma R}\Bigr)e^{-\gamma R/4}+\mt e^{-\lambda_0R}
 +\frac{\mt}{\pi}e^{-\lambda_0R/2},
\end{align*}
using $2\sinh x\le e^x$; and $e^{-\lambda_0R}\le e^{-\lambda_0R/2}$.  The bound by $\mt$ is
Lemma~\ref{lem:transforms}(iv).
\end{proof}

\emph{Comparison with the sketch.}  Step (a) of the sketch states the bound
$\mt[(1+4/(\pi\gamma R))e^{-\gamma R/4}+(1+\lambda_0R/\pi)e^{-\lambda_0R/2}]$.  It also follows
from Step~3, with $J\le x\cosh x\le xe^x$, $x=\lambda_0R/2$, in place of $J\le\sinh x$; for
$\lambda_0R\ge1$ the form of Lemma~\ref{lem:cluster} is the sharper one.

\begin{tlemma}[resolvent form of the state derivative; replaces Lemma 10.1]\label{lem:deriv}
In the path setting of \S8 of the record, assume (G$_\gamma$) and let $\lambda_0>0$.  For
self-adjoint $A\in B(\HH)$ and $0<s<1$ put
\[
 u_{j,s}:=(\lambda_0-iH_s)^{-1}Q_sV_j\Om_s,\qquad
 \Ft_{j,s}(t):=\ip{u_{j,s}}{e^{itH_s}Q_sA\Om_s} .
\]
Then, the sum running over the finitely many $j$ with $h_j\ne0$,
\begin{equation}
 \frac{d}{ds}\,\omega_{g_s}(A)=-2\,\mathrm{Re}\sum_j\Bigl[i\,\Ft_{j,s}(0)
 +\lambda_0\int_\R W_\gamma(t)\,\Ft_{j,s}(t)\,dt\Bigr],
 \label{eq:deriv}
\end{equation}
and the integrals converge absolutely.
\end{tlemma}

\begin{proof}
\emph{Step 1.}  The first display of the proof of Lemma~10.1, which uses Proposition~8.2(ii)
and the self-adjointness of $H_s^{-1}Q_s$, gives
$\frac d{ds}\omega_{g_s}(A)=-2\mathrm{Re}\ip{V(v)\Om_s}{H_s^{-1}Q_sA\Om_s}$, with
$\Om_s\in D(V(v))$ by Proposition~2.6 ($v=g'-g$ is real, smooth, of compact support).

\emph{Step 2.}  The first assertion of Lemma~10.2 (only finitely many $h_j$ are nonzero, and
$V(v)\xi=\sum_jV_j\xi$ on common domains) is proved there from the compact support of $v$ and
the finite additivity of \S2.1 alone; the second assertion of Lemma~10.2, the bound by $M$, is
not used here.  Each $h_j$ is real, bounded, of compact support, so $\Om_s\in D(V_j)$ by
Proposition~2.6; hence $\ip{V(v)\Om_s}{y}=\sum_j\ip{V_j\Om_s}{y}$ for every vector $y$.

\emph{Step 3: a functional-calculus identity.}  On $\spec H_s\subseteq\{0\}\cup[\gamma,\infty)$
the bounded Borel functions $E^{-1}\mathbf 1_{[\gamma,\infty)}(E)$ and
$(\lambda_0+iE)^{-1}\bigl(i+\lambda_0E^{-1}\bigr)\mathbf 1_{[\gamma,\infty)}(E)$ coincide: for
$E\ge\gamma$ the second is $(\lambda_0+iE)^{-1}(iE+\lambda_0)E^{-1}=E^{-1}$, and at $E=0$ both
vanish.  The functional calculus being multiplicative,
$H_s^{-1}Q_s=(\lambda_0+iH_s)^{-1}\bigl(iQ_s+\lambda_0H_s^{-1}Q_s\bigr)$.

\emph{Step 4.}  Fix $j$ and drop the index $s$.  By Step~3,
\[
 \ip{V_j\Om}{H^{-1}QA\Om}=i\ip{V_j\Om}{(\lambda_0+iH)^{-1}QA\Om}
 +\lambda_0\ip{V_j\Om}{H^{-1}Q\,(\lambda_0+iH)^{-1}QA\Om} .
\]
Since $(\lambda_0+iH)^{-1}QA\Om=Q(\lambda_0+iH)^{-1}QA\Om$, the first inner product is
$\ip{(\lambda_0-iH)^{-1}QV_j\Om}{QA\Om}=\Ft_j(0)$.  For the second, Lemma~9.1(ii), whose
hypothesis is (G$_\gamma$) and in which $1-P_{\ker H}=Q$ because $\dim\ker H=1$, gives
$H^{-1}Q=\int W_\gamma(t)e^{itH}Q\,dt$ as a norm-convergent Bochner integral; applied to
$(\lambda_0+iH)^{-1}QA\Om$ and paired with $V_j\Om$,
\begin{align*}
 \ip{V_j\Om}{H^{-1}Q(\lambda_0+iH)^{-1}QA\Om}
 &=\int W_\gamma(t)\ip{(\lambda_0-iH)^{-1}QV_j\Om}{e^{itH}QA\Om}\,dt\\
 &=\int W_\gamma(t)\Ft_j(t)\,dt .
\end{align*}
The integral converges absolutely because $W_\gamma\in L^1$ (Lemma~9.1(i)) and
$|\Ft_j|\le\|u_j\|\,\|A\|$ (Lemma~\ref{lem:transforms}(iv)).  Summing over $j$ (Step~2) and
inserting in Step~1 gives (\ref{eq:deriv}).
\end{proof}

\emph{Domains} (item (b) of the commission).  Every operator applied to $Q_sA\Om_s$ above is
bounded: $\|H_s^{-1}Q_s\|\le\gamma^{-1}$ by (G$_\gamma$), and
$\|(\lambda_0\pm iH_s)^{-1}\|\le\lambda_0^{-1}$.  The only unbounded operator is $V_j$, applied
to $\Om_s\in D(V_j)$.  No domain condition on $Q_sA\Om_s$ is needed; without the gap,
$H_s^{-1}Q_s$ would be unbounded and Step~3 would need $Q_sA\Om_s\in D(H_s^{-1})$.

\begin{tlemma}[two-regime bound with $\Mt$; replaces Lemma 10.3]\label{lem:tworegime}
In the setting of Lemma~\ref{lem:deriv} let $\lambda_0\ge1$, let $O$ be a bounded open interval,
$A=A^*\in\A(O)$ with $\|A\|\le1$, and let $h_j\ne0$ with $r_j=\dist(B_j,O)>0$.  Then
\begin{gather*}
 |\Ft_{j,s}(t)|\le\Mt(\lambda_0)\quad(t\in\R),\qquad
 |\Ft_{j,s}(t)|\le\Mt(\lambda_0)\,\Xi_{\gamma,\lambda_0}(r_j/2)\quad(|t|\le r_j/2),\\
 |\Ft_{j,s}(0)|\le\Mt(\lambda_0)\,\Xi_{\gamma,\lambda_0}(r_j) .
\end{gather*}
\end{tlemma}

\begin{proof}
$h_j$ is real and measurable with $\|h_j\|_\infty\le1$ and $\supp h_j\subseteq[j-1,j+1]$ (proof
of Lemma~10.2), and $g_s\in\G\subset\CH$ (Lemma~3.1); so $\|u_{j,s}\|\le\Mt(\lambda_0)$ by
Definition~\ref{def:Mt}, and the first bound is Lemma~\ref{lem:transforms}(iv).  Let $|t|<r_j$
and $\beta_t:=\beta^{g_s}_t=\mathrm{Ad}\,e^{itH_s}$.  Since $e^{-itH_s}\Om_s=\Om_s$ and $Q_s$
commutes with $e^{itH_s}$, $e^{itH_s}Q_sA\Om_s=Q_s\beta_t(A)\Om_s$, so
$\Ft_{j,s}(t)=\ip{u_{j,s}}{Q_s\beta_t(A)\Om_s}$.  Apply Lemma~\ref{lem:cluster} with $g=g_s$,
$X=V_j$, $B=B_j$, the bounded open interval $O_{|t|}$ and $\beta_t(A)\in\A(O_{|t|})$
(Proposition~2.7).  Its hypotheses: (G$_\gamma$); $V_j$ is affiliated with $\A(O')'$ for every
bounded open interval $O'$ with $\dist(B_j,O')>0$ (\emph{Locality} in \S2.1, as
$B_j=\supp h_j$); $\Om_s\in D(V_j)$ (Proposition~2.6); $R:=\dist(B_j,O_{|t|})\ge r_j-|t|>0$,
because $|y-x|\ge\dist(y,O)-\dist(x,O)>r_j-|t|$ for $y\in B_j$, $x\in O_{|t|}$.  With
$\|\beta_t(A)\|=\|A\|\le1$ it gives $|\Ft_{j,s}(t)|\le\Mt(\lambda_0)\Xi(R)\le
\Mt(\lambda_0)\Xi(r_j-|t|)$, since $\Xi$ is nonincreasing (both summands are).  For
$|t|\le r_j/2$, $r_j-|t|\ge r_j/2$; at $t=0$, $O_0=O$ and $R=r_j$.
\end{proof}

\begin{tprop}[one-path estimate with $\Mt$; replaces Proposition 10.4]\label{prop:onepath}
Assume (G$_\gamma$), let $\lambda_0\ge1$ and $M'\ge\Mt(\lambda_0)$.  If $d\ge1$ and
$g\equiv g'\equiv1$ on $O_d$ ($g,g'\in\G$), then
\[
 \|(\omega_g-\omega_{g'})|_{\A(O)}\|\le12\,M'\sum_{n\ge\lfloor d\rfloor}\Bigl(\Xi_{\gamma,\lambda_0}(n)
 +\lambda_0\|W_\gamma\|_1\,\Xi_{\gamma,\lambda_0}(n/2)+\lambda_0\,T_\gamma(n/2)\Bigr)
 =\Psit_{\gamma,\lambda_0}(d) .
\]
The series converges.
\end{tprop}

\begin{proof}
\emph{Step 1.}  $g\equiv g'\equiv1$ on the open interval $O_d$, so the plateau interiors share
$O_d\ne\varnothing$ and the path setting applies; $v=0$ on $O_d$, so every nonzero $h_j$ has
$B_j\subseteq\supp v\subseteq\R\setminus O_d$ and $r_j\ge d\ge1$ (as in Step~1 of
Proposition~10.4).

\emph{Step 2.}  For $A=A^*\in\A(O)$, $\|A\|\le1$, and $0<s<1$, Lemma~\ref{lem:deriv},
$|\mathrm{Re}\,z|\le|z|$, and Lemma~\ref{lem:tworegime}, with the integral split at
$|t|=r_j/2$, give
\begin{align*}
 \Bigl|\frac d{ds}\omega_{g_s}(A)\Bigr|
 &\le2\sum_j\Bigl(|\Ft_{j,s}(0)|+\lambda_0\!\int|W_\gamma||\Ft_{j,s}|\Bigr)\\
 &\le2M'\sum_j\Bigl(\Xi(r_j)+\lambda_0\|W_\gamma\|_1\Xi(r_j/2)+\lambda_0T_\gamma(r_j/2)\Bigr).
\end{align*}

\emph{Step 3: shell count.}  For every integer $n\ge\lfloor d\rfloor$ at most six indices $j$
have $r_j\in[n,n+1)$.  Write $O=(a_O,b_O)$.  $O_d$ has length $b_O-a_O+2d>2$, and $B_j$ is a
compact subset of $[j-1,j+1]$ disjoint from the open interval $O_d=(a_O-d,b_O+d)$, so each of
its points satisfies $y\le a_O-d$ or $y\ge b_O+d$; points $y\le a_O-d$ and $y'\ge b_O+d$ of $B_j$
would give $y'-y\ge b_O-a_O+2d>2$, contradicting $B_j\subseteq[j-1,j+1]$; so $B_j$ lies on one
side of $O_d$.  On the right side,
$x_j:=\min B_j$ satisfies $x_j-b_O=r_j\in[n,n+1)$ and $|x_j-j|\le1$, so
$j\in[b_O+n-1,b_O+n+2)$, a half-open interval of length $3$, which contains at most three
integers.  On the left side, $x_j:=\max B_j$ gives $j\in(a_O-n-2,a_O-n+1]$, again at most
three.  (Step~3 of Proposition~10.4 says ``at most four integers'' per side; that bound is
true and weaker.)

\emph{Step 4.}  The three functions of $r_j$ in Step~2 are nonincreasing ($\Xi$ by inspection,
$T_\gamma$ by definition), and every $r_j\ge d$ lies in a shell $[n,n+1)$ with
$n\ge\lfloor d\rfloor$; so $|\frac d{ds}\omega_{g_s}(A)|\le12M'\sum_{n\ge\lfloor d\rfloor}(\dots)$.
Integrate over $s\in[\varepsilon,1-\varepsilon]$ and let $\varepsilon\downarrow0$, using the
norm continuity of $s\mapsto P_s$ at the endpoints (Proposition~8.2(i)), as in Step~4 of
Proposition~10.4; $\omega_g-\omega_{g'}$ is hermitian, so its norm on $\A(O)$ is the supremum
over self-adjoint contractions.  Convergence: for $n\ge1$,
$\Xi(n)\le(1+\frac4{\pi\gamma})e^{-\gamma n/4}+(1+\frac1\pi)e^{-\lambda_0n/2}$ and
$\Xi(n/2)\le(1+\frac8{\pi\gamma})e^{-\gamma n/8}+(1+\frac1\pi)e^{-\lambda_0n/4}$, and
$T_\gamma(n/2)\le C_{q,\gamma}(1+n/2)^{-q}$ for every $q\ge0$ (Lemma~9.1(i)).
\end{proof}

\begin{tlemma}[the rate function; replaces Lemma 10.5]\label{lem:rate}
Assume (G$_\gamma$), let $\lambda_0\ge1$ and $\Mt(\lambda_0)\le M'<\infty$, and let
$\Psit:=\Psit_{\gamma,\lambda_0}$ be defined by (\ref{eq:Psit}) and
$\Psit(d):=\max\{2,\Psit(1)\}$ for $0\le d<1$.  Then $\Psit$ is nonincreasing; for every
$N\ge1$ there is $C_N=C_{N,\gamma,\lambda_0,M'}<\infty$ with $\Psit(d)\le C_N(1+d)^{-N}$ for all
$d\ge0$; and $\|(\omega_g-\omega_{g'})|_{\A(O)}\|\le\Psit(d)$ whenever $g,g'\in\G$ satisfy
$g\equiv g'\equiv1$ on $O_d$, $d\ge0$.
\end{tlemma}

\begin{proof}
Monotonicity: for $d\ge1$, a larger $\lfloor d\rfloor$ removes nonnegative terms; on $[0,1)$
the value $\max\{2,\Psit(1)\}$ dominates every value on $[1,\infty)$.  Decay: for $d\ge1$,
$\lfloor d\rfloor\ge(1+d)/4$ (for $1\le d<2$ because $d<3$; for $d\ge2$ because
$\lfloor d\rfloor\ge d-1\ge(1+d)/4$).  For $\beta>0$,
$\sum_{n\ge\lfloor d\rfloor}e^{-\beta n}=e^{-\beta\lfloor d\rfloor}(1-e^{-\beta})^{-1}\le
e^{-\beta(1+d)/4}(1-e^{-\beta})^{-1}$, which is $\le C(1+d)^{-N}$; this covers the
$\Xi$-terms, by the bounds at the end of the proof of Proposition~\ref{prop:onepath}.  For the
$T_\gamma$-terms take $q:=N+1$: $T_\gamma(n/2)\le C_{q,\gamma}2^q(1+n)^{-q}$ and
$\sum_{n\ge\lfloor d\rfloor}(1+n)^{-q}\le\int_{\lfloor d\rfloor-1}^\infty(1+x)^{-q}dx=
\lfloor d\rfloor^{1-q}/(q-1)\le4^{q-1}(1+d)^{1-q}/(q-1)$.  For $0\le d<1$,
$\Psit(d)\le2^N\max\{2,\Psit(1)\}(1+d)^{-N}$.  The last assertion is
Proposition~\ref{prop:onepath} for $d\ge1$ and the trivial bound
$\|\omega_g-\omega_{g'}\|\le2$ for $d<1$.
\end{proof}

\section{The Euclidean side: steps (c) and (d) of the sketch}\label{sec:euclidean}

\begin{tlemma}[the elementary inequality]\label{lem:ineq}
The function $k$ of (\ref{eq:k}) satisfies $k(E)=2\int_0^1(1-u)e^{-uE}du$; it is continuous,
positive and nonincreasing on $[0,\infty)$, and $k(2)=(1+e^{-2})/2=1/c_k$.  For $\lambda_0\ge1$
and $E\ge0$,
\[
 (\lambda_0^2+E^2)^{-1}\le c_k\,k(E) .
\]
\end{tlemma}

\begin{proof}
For measurable $\varphi\ge0$, $\iint_{[-\frac12,\frac12]^2}\varphi(\tau-\tau')\,d\tau\,d\tau'
=\int_{-1}^1(1-|u|)\varphi(u)\,du$, because for $|u|\le1$ the set of $\tau$ with
$\tau,\tau-u\in[-\frac12,\frac12]$ has length $1-|u|$.  With $\varphi(u)=e^{-|u|E}$ this gives
$k(E)=2\int_0^1(1-u)e^{-uE}du$; at $E=0$ the value is $1$, and for $E>0$ an integration by parts
gives $\int_0^1(1-u)e^{-uE}du=E^{-1}-(1-e^{-E})E^{-2}=(E-1+e^{-E})E^{-2}$, which is
(\ref{eq:k}).  The integrand is positive and nonincreasing in $E$ for each $u\in[0,1)$, so $k$
is positive and nonincreasing, and continuous by dominated convergence;
$k(2)=2(1+e^{-2})/4$.  \emph{Case $0\le E\le2$:} $c_kk(E)\ge c_kk(2)=1\ge\lambda_0^{-2}\ge
(\lambda_0^2+E^2)^{-1}$.  \emph{Case $E\ge2$:} $k(E)\ge2(E-1)E^{-2}\ge E^{-1}$ (the second
inequality is $E\ge2$), $E^{-1}\ge(1+E^2)^{-1}$ (as $E^2-E+1>0$), $(1+E^2)^{-1}\ge
(\lambda_0^2+E^2)^{-1}$ (as $\lambda_0\ge1$), and $c_k\ge1$.
\end{proof}

\emph{Tests.}  By hand: $E=0$ gives $\lambda_0^{-2}\le1<c_k$; $E=2$, $\lambda_0=1$ gives
$\frac15\le1$; $E=10$ gives $\frac1{101}\le c_k\cdot\frac{2(9+e^{-10})}{100}=0.317\ldots$.  The
closed form of $k$ and the inequality on a grid are checked in Appendix~\ref{app:numerics}
(\emph{numerical}).  The constant is not sharp; it is chosen so that the case split falls at
$E=2$.

The next lemma is the Fubini step on which both the time-integrated observable of Route T and
\S\ref{sec:N39} rest.  It concerns the objects of Definition~N.26 (\texttt{aux:d1-obs},
p.~81) of the record: for $\kappa\ge2$ the regularized field
$\Phi_\kappa(\zeta)=\Phi(\delta_\kappa(\zeta-\cdot))$, its variance $c_\kappa$, the Wick power
$A_\kappa(z):={:}\Phi_\kappa(z)^n{:}_{c_\kappa}$, and the regularized monomials, integrated
pointwise in $q\in\mathcal S'(\R^2)$; the localized monomial is their $L^2(\dmu)$-limit
(Lemma~N.28(a), \texttt{aux:d1-moments}).  The squares are
$\Delta_{l,k}=[l-\frac12,l+\frac12)\times[k-\frac12,k+\frac12)$, time first (Definition~N.1).

\begin{tlemma}[time slicing of a one-square time-integrated monomial]\label{lem:slice}
Let $\Delta=\Delta_{l,k}$, let $w:\R^2\to\R$ be bounded, Borel measurable and zero outside
$\Delta$, and let $0\le n\le2p'$.  Let $Q$ be the localized monomial of type (1.1.6) of
Definition~N.26(a) with $r=1$, degree $n$ and kernel $w$, and, for $t\in\R$, let $Q(t)$ be the
localized monomial of type (1.1.5) of Definition~N.26(b) with $r=1$, degree $n$, time $t$ and
kernel $w(t,\cdot)$ (so $Q(t)=0$ for $t\notin[l-\frac12,l+\frac12)$).
\begin{enumerate}
\item[\rm(i)] $\sup_t\|Q(t)\|_{L^2(\dmu)}<\infty$; $t\mapsto Q(t)$ is strongly measurable from
$\R$ to $L^2(\dmu)$; and $Q=\int_\R Q(t)\,dt$ as an $L^2(\dmu)$-valued Bochner integral.  Hence
$\ell(Q)=\int\ell(Q(t))\,dt$ for every bounded linear functional $\ell$ on $L^2(\dmu)$.
\item[\rm(ii)] Let $n\ge1$.  For $t\in[l-\frac12,l+\frac12)$, $Q(t)$ is localized in $\Delta$
(Remark~N.2(ii): the row of the time $t$ is $l$), with local degree $N(\Delta)=n$ and
$\GJS{Q(t)}=K_1K_2^{2n}n!\,\|w(t,\cdot)\|_{L^2(\R)}$; and
$\int_\R\GJS{Q(t)}\,dt\le K_1K_2^{2n}n!\,\|w\|_{L^2(\R^2)}=\GJS{Q}$.
\end{enumerate}
\end{tlemma}

\begin{proof}
Write $Q^{(\kappa)}(q):=\int w(z)A_\kappa(z)(q)\,dz$ and
$Q^{(\kappa)}(t)(q):=\int w(t,x)A_\kappa(t,x)(q)\,dx$ (Definition~N.26).

\emph{Step 1 (Fubini, pointwise).}  For each $q$, $z\mapsto A_\kappa(z)(q)$ is continuous
(Definition~N.26: $\zeta\mapsto\Phi_\kappa(\zeta)(q)$ is smooth, and (35) is a polynomial), and
$w$ is bounded with compact support; Fubini on $\R^2$ gives
$Q^{(\kappa)}(q)=\int Q^{(\kappa)}(t)(q)\,dt$ for every $q$.

\emph{Step 2 (moments, uniformly in $t$ and $\kappa$).}  Lemma~N.24 with $\ell=2$ and the
Fubini step of the proof of Lemma~N.28 (display (37) with $j=1$, and the majorant (38) for type
(b)) give
$\E[Q^{(\kappa)}(t)^2]=n!\iint w(t,x)w(t,y)G_{\kappa\kappa}((0,x-y))^n\,dx\,dy
\le n!\,\|w\|_\infty^2\iint_{[k-\frac12,k+\frac12]^2}M_1(|x-y|)^n\,dx\,dy=:c_w$,
finite by Lemma~N.25(iii)--(iv) and independent of $t$ and of $\kappa\ge2$.  Also
$\E[A_\kappa(z)^2]=n!\,c_\kappa^n$ (Lemma~N.24, $\ell=2$, equal variables).

\emph{Step 3 (the regularized slices form a Bochner integral).}  Let $\phi\in L^2(\dmu)$.  The
function $(t,x,q)\mapsto\phi(q)w(t,x)A_\kappa(t,x)(q)$ is measurable ($A_\kappa(z)(q)$ is
continuous in $z$ and measurable in $q$, hence jointly measurable), and
$\iiint|\phi\,w\,A_\kappa|\,d\mu_0\,dx\,dt\le\|\phi\|_2(n!c_\kappa^n)^{1/2}\|w\|_{L^1(\R^2)}<\infty$.
By Fubini, $t\mapsto\ip{\phi}{Q^{(\kappa)}(t)}$ is measurable and, with Step~1,
$\int\ip{\phi}{Q^{(\kappa)}(t)}dt=\ip{\phi}{Q^{(\kappa)}}$.  $L^2(\dmu)$ is separable: its
$\sigma$-algebra is generated by $\Phi(f_m)$, $\{f_m\}$ a countable dense subset of
$\mathcal S(\R^2)$, because $q(f)=\lim q(f_m)$ when $f_m\to f$ in $\mathcal S(\R^2)$.  By
Pettis's theorem $t\mapsto Q^{(\kappa)}(t)$ is strongly measurable; it is bounded by
$c_w^{1/2}$ (Step~2) and vanishes outside $[l-\frac12,l+\frac12)$, hence Bochner integrable,
and $\ip{\phi}{\int Q^{(\kappa)}(t)dt}=\ip{\phi}{Q^{(\kappa)}}$ for all $\phi$, i.e.\
$\int Q^{(\kappa)}(t)\,dt=Q^{(\kappa)}$.

\emph{Step 4 ($\kappa\to\infty$).}  By Lemma~N.28(a), applied to type (a) and to type (b),
$Q^{(\kappa)}\to Q$ and $Q^{(\kappa)}(t)\to Q(t)$ in $L^2(\dmu)$ for every $t$.  Fix
$\kappa_m\to\infty$.  Then $t\mapsto Q(t)$ is strongly measurable, as a pointwise limit of
strongly measurable functions, and $\|Q(t)\|_2^2\le c_w$; moreover
$\|Q^{(\kappa_m)}(t)-Q(t)\|_2\le2c_w^{1/2}$, it tends to $0$ for every $t$, and it vanishes
outside $[l-\frac12,l+\frac12)$.  By dominated convergence
$\int\|Q^{(\kappa_m)}(t)-Q(t)\|_2\,dt\to0$, so
$\int Q(t)\,dt=\lim_m\int Q^{(\kappa_m)}(t)\,dt=\lim_mQ^{(\kappa_m)}=Q$.  Bounded linear
functionals commute with Bochner integrals.

(ii) For $t\in[l-\frac12,l+\frac12)$, $w(t,x)\ne0$ only if $(t,x)\in\Delta$, i.e.\
$x\in[k-\frac12,k+\frac12)$; so $Q(t)$ is localized in $\Delta_{l,k}$ (Remark~N.2(ii)), with
$N(\Delta)=n$ and all other local degrees $0$, and (1.1.8) gives its norm.  Cauchy--Schwarz on
an interval of length $1$ gives $\int_{l-1/2}^{l+1/2}\|w(t,\cdot)\|_2\,dt\le
(\int\|w(t,\cdot)\|_2^2dt)^{1/2}=\|w\|_{L^2(\R^2)}$, and (1.1.7) gives
$\GJS{Q}=K_1K_2^{2n}n!\|w\|_{L^2(\R^2)}$ (as in Lemma~N.18(i) for $U_{\Delta,n}$).
\end{proof}

\emph{Borel versions.}  Every kernel used below is bounded, measurable and of compact support.
Replacing it by a Borel function equal to it almost everywhere changes neither the regularized
integrals (pointwise in $q$) nor their limits, and changes the slices only for a null set of
times; so the Borel hypothesis of Lemma~\ref{lem:slice} is no restriction.

\begin{tlemma}[the time-integrated interaction]\label{lem:Vbar}
Let $h'$ be real and measurable with $\|h'\|_\infty\le1$ and $\supp h'\subseteq[a,a+2]$, and let
$k_a$, $S_a=\{\Delta_{0,k}:k_a\le k\le k_a+2\}$ be as in Definition~N.10
(\texttt{aux:notation-ringdef}).  For $k\in\{k_a,k_a+1,k_a+2\}$ and $0\le n\le2p'$ let
$h'_k:=h'\mathbf 1_{[k-\frac12,k+\frac12)}$, $w_{n,k}(\tau,x):=\lambda a_n\mathbf
1_{[-\frac12,\frac12)}(\tau)\,h'_k(x)$, let $\Vb_{n,k}$ be the localized monomial of type (1.1.6)
with $r=1$, degree $n$ and kernel $w_{n,k}$, and put $\Vb(h'):=\sum_{n,k}\Vb_{n,k}$.
\begin{enumerate}
\item[\rm(i)] $\tau\mapsto J_\tau V(h')$ is strongly measurable from $[-\frac12,\frac12)$ to
$L^2(\dmu)$, and $\Vb(h')=\int_{-1/2}^{1/2}J_\tau V(h')\,d\tau$ as an $L^2(\dmu)$-valued Bochner
integral.  Here $V(h')$ is the $V(h')$ of \S2.1 of the record and $J_\tau$ the sharp-time map of
the proof of Lemma~N.22 (\texttt{aux:notation-sharpid}).
\item[\rm(ii)] $\Vb_{n,k}={:}\Phi^n(w_{n,k}){:}\in\Gamma_n(\HH_G)$, so
$\Vb(h')\in\bigoplus_{n=0}^{2p'}\Gamma_n(\HH_G)$;
$\|\Vb(h')\|_{L^2(\dmu)}\le\|V(h')\|_{L^2(\dphi)}\le V_*^{1/2}$;
$\|\Vb(h')\|_{L^4(\dmu)}\le3^{p'}\|\Vb(h')\|_{L^2(\dmu)}$; hence
$\|\Vb(h')^2\|_{L^2(\dmu)}\le9^{p'}V_*$.
\item[\rm(iii)] For $0\le n,n'\le2p'$ and $k,k'\in\{k_a,k_a+1,k_a+2\}$ let $\overline
Q_{n,n',k,k'}$ be the localized monomial of type (1.1.6) with $r=2$, degrees $(n,n')$ and
kernel $w_{n,k}\otimes w_{n',k'}\in L^2(\R^4)$, localized in $\Delta_{0,k}$ and $\Delta_{0,k'}$,
which belong to $S_a$.  Then $\overline Q_{n,n',k,k'}=\Vb_{n,k}\Vb_{n',k'}$ almost everywhere,
and $\Vb(h')^2=\sum_{n,n',k,k'}\overline Q_{n,n',k,k'}$ almost everywhere.
\item[\rm(iv)] For this decomposition,
$\sum_{n,n',k,k'}\GJS{\overline Q_{n,n',k,k'}}\le9K_1K_2^{8p'}(4p')!\,\lambda^2(\sum_{n=0}^{2p'}|a_n|)^2$.
\end{enumerate}
\end{tlemma}

\begin{proof}
(i) By Lemma~\ref{lem:slice} with $\Delta=\Delta_{0,k}$ and $w=w_{n,k}$ (a Borel version),
$\Vb_{n,k}=\int\Vb_{n,k}(\tau)\,d\tau$, where for $\tau\in[-\frac12,\frac12)$ the slice
$\Vb_{n,k}(\tau)$ is the monomial of type (1.1.5) with $r=1$, degree $n$, time $\tau$ and
kernel $\lambda a_nh'_k$, and $\Vb_{n,k}(\tau)=0$ otherwise.  By Lemma~N.28(f)(2) (a real
bounded measurable kernel vanishing outside one column, $0\le n\le2p'$),
$\Vb_{n,k}(\tau)=\lambda a_n\,J_\tau W_n(h'_k)$ almost everywhere, for every
$\tau\in[-\frac12,\frac12)$.  By Remark~N.23(b), $\sum_kW_n(h'_k)=W_n(h')$ (linearity in the
kernel), and by Lemma~N.23$'$(iii), $\lambda\sum_na_nW_n(h')=V(h')$; $J_\tau$ is linear.  Hence
$\sum_{n,k}\Vb_{n,k}(\tau)=J_\tau V(h')$ for $\tau\in[-\frac12,\frac12)$; the strong
measurability and the integral formula follow by summing the finitely many Bochner integrals.

(ii) $\Vb_{n,k}$ is the $L^2(\dmu)$-limit of $\int w_{n,k}(z){:}\Phi_\kappa(z)^n{:}_{c_\kappa}dz$
(Definition~N.26(a), $r=1$); Lemma~N.23$'$(ii), in the case $(\HH_G,\Phi,G,2)$ with $\nu=\delta$
and the real bounded kernel $w_{n,k}$ of compact support, identifies the limit with
${:}\Phi^n(w_{n,k}){:}\in\Gamma_n(\HH_G)$.  By (i) and the norm inequality for Bochner integrals,
$\|\Vb(h')\|_2\le\int_{-1/2}^{1/2}\|J_\tau V(h')\|_{L^2(\dmu)}d\tau=\|V(h')\|_{L^2(\dphi)}$, since
$J_\tau$ is isometric (proof of Lemma~N.22), and $\|V(h')\|^2_{L^2(\dphi)}\le V_*$ by
Lemma~N.20(iv) with Lemma~N.23$'$(iii).  The proof of Lemma~N.21(ii) uses of its $Y$ only that
it is a sum of Wick chaoses of degree $\le2p'$ in $L^2(\dmu)$ (it says so); it applies verbatim
to $\Vb(h')$ and gives $\|\Vb(h')\|_4\le3^{p'}\|\Vb(h')\|_2$ (this is also Simon's
Theorem~I.22 at $p=4$, $m=2p'$, as quoted in Assumption~A4).  Finally
$\|\Vb^2\|_2=\|\Vb\|_4^2\le9^{p'}\|\Vb\|_2^2\le9^{p'}V_*$.

(iii) The kernel $w_{n,k}\otimes w_{n',k'}$ is the product of the bounded functions
$w_{n,k}$, $w_{n',k'}$, which vanish outside $\Delta_{0,k}$ and $\Delta_{0,k'}$; by
Lemma~N.28(c) (type (a)), $\overline Q_{n,n',k,k'}=\Vb_{n,k}\Vb_{n',k'}$ almost everywhere, and
$\sum_{n,n',k,k'}\Vb_{n,k}\Vb_{n',k'}=(\sum_{n,k}\Vb_{n,k})^2$.  The individual degrees are
$\le2p'<\bar n=2p'+1$; a factor of degree $0$ contributes $0$ to the local degree, as in
Lemma~N.43(ii).

(iv) The computation of Lemma~N.43(iii), with
\[
 \|w_{n,k}\|_{L^2(\R^2)}=\lambda|a_n|\,\|\mathbf 1_{[-\frac12,\frac12)}\|_{L^2(\R)}\,
 \|h'_k\|_{L^2(\R)}\le\lambda|a_n|
\]
in place of $\lambda|a_n|\,\|h'_k\|_2$: for one ordered column pair the weight of (1.1.7) is
$K_1K_2^{2(n+n')}(n+n')!$ if $k=k'$ and $K_1K_2^{2n}n!\,K_2^{2n'}n'!$ if $k\ne k'$, in both cases
$\le K_1K_2^{8p'}(4p')!$ ($K_2\ge1$, $n+n'\le4p'$, $n!n'!\le(n+n')!$); the kernel norm is
$\le\lambda^2|a_na_{n'}|$; there are at most $9$ ordered column pairs; and
$\sum_{n,n'}|a_na_{n'}|=(\sum_n|a_n|)^2$.
\end{proof}

\begin{ttheorem}[time-integrated copy of Theorem D.4]\label{thm:TD4}
Assume the standing assumption of \S D.c.1 of the record (\texttt{sub:d2-conv}).  For every
$h\in\CE$, every $a\in\R$ and every real measurable $h'$ with $\|h'\|_\infty\le1$ and
$\supp h'\subseteq[a,a+2]$,
\[
 \int\Vb(h')^2\,dq_h\le M_B^2 ,
\]
with $M_B^2=e^{15\lambda|a_0|}S_2(B_{3\times5})\,9^{p'}V_*+C_7c_U\overline\Sigma_{3\times5}\cdot
9K_1K_2^{8p'}(4p')!\,\lambda^2(\sum_n|a_n|)^2$, the constant of Theorem~D.4.
\end{ttheorem}

\begin{proof}
\emph{Step 1.}  By Lemma~\ref{lem:Vbar}(iii), $\Vb(h')^2$ is almost everywhere the finite sum
of the localized monomials $\overline Q_{n,n',k,k'}$ of type (1.1.6), with kernels in $L^2$,
individual degrees $<\bar n$, and localization squares in $R:=S_a$; $B:=B_{3\times5}(a)$ is a
box containing $R$.  $\Vb(h')\in L^4(\dmu)$, so $\Vb(h')\in L^2(dq_h)$ for every $h\in\CE$
(Lemma~N.32, \texttt{aux:d2-int}) and $\Vb(h')^2\in L^1(dq_h)$.  Since $\Vb^2\ge0$,
$\int\Vb^2dq_h=\ev{\Vb^2}{h^B}+(\ev{\Vb^2}{h}-\ev{\Vb^2}{h^B})$ with $\ev{\Vb^2}{h^B}\ge0$.

\emph{Step 2 (base term).}  As in Step~2 of Theorem~D.4: $h^B\in\CE$ with $\supp
h^B\subseteq\overline B$, of measure $15$ (Lemma~N.5(i)); Lemma~C.2 (\texttt{lem:Z}) gives
$Z(h^B)\ge e^{-15\lambda|a_0|}$; Proposition~N.15(iv),(iii) gives
$\|e^{-\Uc_{h^B}}\|_{L^2(\dmu)}\le S_2(B_{3\times5})$.  By Cauchy--Schwarz and
Lemma~\ref{lem:Vbar}(ii),
$\ev{\Vb^2}{h^B}\le e^{15\lambda|a_0|}\|\Vb^2\|_{L^2(\dmu)}S_2(B_{3\times5})\le
e^{15\lambda|a_0|}9^{p'}V_*S_2(B_{3\times5})$.

\emph{Step 3 (strips).}  Exactly Step~3 of Theorem~D.4, which depends only on $R=S_a$ and $B$:
each $\Delta=\Delta_{l,k}\not\subseteq B$ lies in one ring $\mathcal L_j$, $j\ge1$, and receives
a strip of width $j$, parallel to the time-zero line if $|l|=1+j$ (class (P)) and bounded by two
lines of constant $x^1$ if $|l|\le j$ (class (T)) (Lemma~N.11(iv)); hence
$\sum_{\Delta\in\mathcal D(h,B)}e^{-m_1w_\Delta}\le\overline\Sigma_{3\times5}$
(Lemma~N.11(ii),(vi)).

\emph{Step 4 (comparison term).}  Lemma~D.2(a) (\texttt{lem:D2}), with $F=\sum\overline
Q_{n,n',k,k'}$ and these pieces, $R$, $B$, $h$ and the strips of Step~3, together with
Lemma~\ref{lem:Vbar}(iv), gives
\[
 |\ev{\Vb^2}{h}-\ev{\Vb^2}{h^B}|\le C_7c_U\GJS{\Vb^2}\,\overline\Sigma_{3\times5}
 \le C_7c_U\overline\Sigma_{3\times5}\cdot9K_1K_2^{8p'}(4p')!\,\lambda^2\Bigl(\sum_n|a_n|\Bigr)^2 .
\]
The hypotheses of Lemma~D.2 hold: $F$ is a finite sum of localized monomials of type (1.1.6),
which the lemma admits (``of type (1.1.5) or (1.1.6)''); $R\subseteq B$; $h\in\CE$; every
strip has positive width $j\ge1$ and separates $\Delta$ from $\bigcup S_a$.  The estimates
inside Lemma~D.2 are Source Fact~2.2 (\texttt{sf:CE2}; Assumption~A8(iv)) for the pairs
\begin{itemize}
\item $Q_1=\overline Q_{n,n',k,k'}$: type (1.1.6), $r=2$, degrees $n,n'\le2p'<\bar n$,
localized in $\Delta_{0,k},\Delta_{0,k'}\in S_a$, kernel $w_{n,k}\otimes w_{n',k'}\in L^2(\R^4)$;
\item $Q_2=U_{\Delta,n''}$, $1\le n''\le2p'$, $\Delta\in\mathcal D(h,B)\cap\mathcal L_j$: type
(1.1.6), $r=1$, localized in $\Delta$, kernel $\lambda a_{n''}h\mathbf 1_\Delta$
(Lemma~N.18(i));
\item the cutoff $h_\alpha=h^B+\alpha(h-h^B)\in\CE$, $0\le\alpha\le1$ (Lemma~N.5(ii));
\item the strip of Step~3, of width $j$: parallel (class (P), case P3 of \S\ref{sec:cases}) or
transverse, bounded by two lines of constant $x^1$ (class (T), case P4).
\end{itemize}
Because $Q_1$ is of type (1.1.6), none of these applications is the case A8(iv)-T as the record
defines it (Assumption~A8, item (iv-T): ``Source Fact~2.2 applied to a sharp-time monomial
$Q_1$ of type (1.1.5)''); all are applications of A8(iv) at its printed scope.

\emph{Step 5.}  Steps~1, 2 and 4 give $\int\Vb^2dq_h\le M_B^2$.  No constant depends on $h$,
$a$ or $h'$ (Step~5 of Theorem~D.4).
\end{proof}

\emph{Status of Theorem~\ref{thm:TD4}.}  Proved, conditional on A8(iv) at its printed scope in
the cases P3--P4, and under the reading of Remark~N.2(iv) (\texttt{aux:notation-gjslattice};
A8 Scope (d)) that the constructed objects, here the $\overline Q_{n,n',k,k'}$ and the
$U_{\Delta,n''}$, are GJS's monomials of p.~593 and that $\Uc_h$ is GJS's $V(h)$.  All other
inputs are statements of the record, used as it states them: Lemmas D.1, D.2, C.2,
Proposition~N.15, Lemmas N.5, N.11, N.18, N.20--N.23$'$, N.24--N.28, N.32, and Assumptions A3,
A4.

\subsection{Theorem D.4 step by step, for $\Vb(h')^2$ in place of $V(h')(0)^2$}

Item (d) of the commission asks, for each step of the proof of Theorem~D.4 and of what it
invokes, whether it holds unchanged, holds with a stated change, or fails.  None fails.

\begin{center}\small
\begin{tabular}{@{}p{0.43\textwidth}p{0.51\textwidth}@{}}
\toprule
item of Theorem D.4 (record) & for $\Vb(h')^2$\\
\midrule
setting: $V(h')$ constructed (N.20), $Y=J_0V(h')$, $Q_V=Y^2$ & changed:
$\Vb(h')=\sum\Vb_{n,k}=\int_{-1/2}^{1/2}J_\tau V(h')d\tau$ (Lemmas~\ref{lem:slice},
\ref{lem:Vbar}(i)), and $\Vb(h')^2$\\
box $B_{3\times5}(a)$, $R=S_a$, rings $\mathcal L_j$, ring constant $\overline\Sigma_{3\times5}$ &
unchanged: the time support $[-\frac12,\frac12)$ of $\Vb(h')$ is row $0$, and its columns are
those of $h'$\\
$V_*$ & unchanged; used through $\|\Vb\|_2\le\|V(h')\|_2$ (Bochner norm inequality, $J_\tau$
isometric) instead of $\|Y\|_2=\|V(h')\|_2$\\
Lemma N.43(i): $Q_V$ as a sum of pieces (N.22, N.23(b)) & changed: Lemma~\ref{lem:Vbar}(iii),
through N.28(c) for type (a)\\
Lemma N.43(ii): pieces of type (1.1.5), two factors at time $0$ & changed: pieces of type
(1.1.6), $r=2$, two factors integrated over row $0$\\
Lemma N.43(iii): $\GJS{Q_V}\le9K_1K_2^{8p'}(4p')!\lambda^2(\sum|a_n|)^2$ & same bound; each kernel
norm acquires the factor $\|\mathbf 1_{[-\frac12,\frac12)}\|_2=1$\\
Step 1 (positivity, splitting) & unchanged\\
Step 2: Lemma C.2, $S_2(B_{3\times5})$, Proposition N.15 & unchanged (the measures are the same)\\
Step 2: $\|Y\|_4\le3^{p'}\|Y\|_2$ (N.21(ii)) & holds with a stated change: applied to
$\Vb(h')\in\bigoplus_{n\le2p'}\Gamma_n(\HH_G)$ (N.23$'$(ii)); the argument of N.21(ii) uses
nothing else\\
Step 3 (strips (P) and (T), $\overline\Sigma_{3\times5}$) & unchanged\\
Step 4: Lemma D.2(a), with Lemma D.1 (class $\mathfrak F$ contains type (1.1.6)) and N.32 &
unchanged in form; the first monomial of each pair is now of type (1.1.6) with $r=2$\\
Step 4: ``every application to a square of class (T) is \dots A8(iv)-T'' & changed: class (T)
is now the case P4 (two time-integrated monomials), not A8(iv)-T\\
Step 5 (constants) & unchanged: the constant is $M_B^2$\\
Appendix C (C.2, C.4 through Step 3 of D.2, C.5 through $S_2$) & unchanged: they concern the
measures $dq_{h_\alpha}$ and $\Uc_{h-h^B}$, which do not change\\
\bottomrule
\end{tabular}
\end{center}

\begin{tlemma}[from the Euclidean bound to the resolvent-weighted quantity]\label{lem:fkn}
Let $g\in\CH$, $H=H(g)$, $\Om=\Om_g$, $c_g=\ip{\Om_g}{\Om_0}>0$ and
$\psi_s=e^{-sH}\Om_0$ (\S4 of the record); let $h$ be real and measurable with
$\|h\|_\infty\le1$ and $\supp h\subseteq[a,a+2]$, and $X:=V(h)$.
\begin{enumerate}
\item[\rm(i)] For every $s>0$, $Xe^{-sH}$ is a bounded operator.
\item[\rm(ii)] For $T\ge1$ and $-\frac12\le\tau<\tau'\le\frac12$,
\[
 \int J_\tau X\cdot J_{\tau'}X\,dq_{h_{T,g}}
 =\frac{\ip{X\psi_{T+\tau}}{e^{-(\tau'-\tau)H}X\psi_{T-\tau'}}}{\|\psi_T\|^2} .
\]
\item[\rm(iii)] $\lim_{T\to\infty}\int\Vb(h)^2\,dq_{h_{T,g}}=\ip{X\Om}{k(H)X\Om}$.
\end{enumerate}
\end{tlemma}

\begin{proof}
(i) By Assumption~A5(i) (\texttt{A:domain}), $\Ran e^{-sK(g)}\subseteq D(X)$, and
$e^{-sH}=e^{sE(g)}e^{-sK(g)}$ has the same range; so $Xe^{-sH}$ is everywhere defined.  It is
closed: if $\xi_m\to\xi$ and $Xe^{-sH}\xi_m\to\eta$, then $e^{-sH}\xi_m\to e^{-sH}\xi$, and $X$
is closed, so $Xe^{-sH}\xi=\eta$.  The closed graph theorem makes it bounded.

(ii) \emph{Step 0: clipping commutes with $J_\tau$.}  For $\ell>0$ let
$c_\ell(y):=(-\ell)\vee(y\wedge\ell)$; it is $1$-Lipschitz.  For a cylinder function $F$ the
identity $J_\tau(c_\ell\circ F)=c_\ell\circ J_\tau F$ is the definition of $J_\tau$ (proof of
Lemma~N.22).  For general $F\in L^2(\dphi)$ approximate $F$ in $L^2$ by cylinder functions, the
density by which that proof extends $J_\tau$ to $L^p$: $c_\ell\circ(\cdot)$ is continuous in
$L^2$, and $J_\tau$ is an $L^2$-isometry.  So $J_\tau(c_\ell\circ F)=c_\ell\circ J_\tau F$.  In
particular $|J_\tau F|\le C$ almost everywhere whenever $|F|\le C$, since then
$J_\tau F=J_\tau(c_C\circ F)=c_C\circ J_\tau F$.

\emph{Step 1: bounded functions.}  Let $F_1,F_2$ be real measurable functions of the time-zero
field with $|F_i|\le C$, and let $-T<\tau<\tau'<T$.  We show
\begin{equation}
 \ip{\Om_0}{e^{-(T-\tau')K}F_2\,e^{-(\tau'-\tau)K}F_1\,e^{-(T+\tau)K}\Om_0}
 =\int J_\tau F_1\cdot J_{\tau'}F_2\;e^{-\Uc_{h_{T,g}}}\,d\mu_0 ,\qquad K=K(g).
 \label{eq:fknasym}
\end{equation}
This is an instance of Assumption~A2 (\texttt{A:fkn}, ll.~825--838), stated ``For all bounded
time-zero functions $f,h'$, $T>0$, and $g\in\CH$'', ``together with the analogous multi-time
formulas on finite slabs'': the functions $1,F_1,F_2,1$ sit at the times $-T,\tau,\tau',T$ of the
slab $[-T,T]$, and A2's display puts the later time on the left.  The record uses these formulas
in Lemma~4.1 at the symmetric times $\mp t/2$ and, in its formula (3), at one time $s$ placed
anywhere in $[-T,T]$.  A2's wording does not restrict the positions of the times.

The printed theorems behind A2, however, have functions of finitely many fields at the interior
times.  Simon's Theorem~V.15 (p.~163) is printed for ``measurable functions on $R^k$ with
$|G_j(x)|\le C(1+|x|^N)$'' of the fields $\phi_F(f_1),\dots,\phi_F(f_k)$, and GRS's
Theorem~II.16 likewise.  For other functions Simon has only the Remark on p.~164, printed without
proof.  So (\ref{eq:fknasym}) is derived here from the bounded case of Theorem~V.15, whose proof
is printed (``Suppose first that $G_0,\dots,G_n$ are bounded \dots\ This proves (V.47) in case are
bounded'', p.~164), by approximation.
\begin{enumerate}
\item[(1)] \emph{Approximants.}  Polynomials in finitely many fields $\varphi_0(f)$, $f$ real
Schwartz functions, are dense in $L^2(\dphi)$.  This follows from Assumption~A1(c),
$L^2(\dphi)=\Gamma(\HH_C)$, the closed span of the Wick monomials, together with the density of
$\mathcal S(\R)$ in $\HH_C$; it is also the density by which the proof of Lemma~N.22 extends
$J_\tau$.  Choose such polynomials $P_{i,m}\to F_i$ in $L^2(\dphi)$ and put
$F_i^{(m)}:=c_C\circ P_{i,m}$.  Then $F_i^{(m)}$ is a bounded function of finitely many fields
with $|F_i^{(m)}|\le C$.  Since $c_C$ is $1$-Lipschitz and $c_C\circ F_i=F_i$,
$\|F_i^{(m)}-F_i\|_2\le\|P_{i,m}-F_i\|_2\to0$.
\item[(2)] \emph{(\ref{eq:fknasym}) for the approximants.}  Apply Theorem~V.15 with $n=3$,
$s_0=-T$, $t_1=T+\tau$, $t_2=\tau'-\tau$, $t_3=T-\tau'$ (all positive), $G_0=G_3=1$, and $G_1$,
$G_2$ the bounded functions that define $F_1^{(m)}$, $F_2^{(m)}$.  The fields of both
approximants are listed among $f_1,\dots,f_k$.  The hypotheses hold: a bounded function satisfies
$|G_j(x)|\le C(1+|x|^N)$ with $N=0$; and $g\ge0$ is bounded with compact support, so
$g\in L^1\cap L^{1+\epsilon}$.  Then $s_1=\tau$, $s_2=\tau'$, $s_3=T$.  The dictionary is A2's:
Simon's $H=H_0+H_I(g)$ is the record's $K(g)$, his $U(g\otimes(s_0,s_3))$ is $\Uc_{h_{T,g}}$,
his $J_s$ is the record's $J_s$, and his $\phi_F$ is the record's $\varphi_0$.  Theorem~V.15
puts the earliest time on the left.  All the operators are self-adjoint and the right side is
real, so taking the complex conjugate of its left side gives the ordering of
(\ref{eq:fknasym}).  That the real Schwartz functions belong to Simon's class $F$ (``Let
$f_1,\dots,f_k\in F$'', p.~163, where $F$ is not restated) is the reading under which A2 cites
Theorem~V.15 for the record's time-zero fields.
\item[(3)] \emph{The Hilbert side converges.}  Let $y:=e^{-(T+\tau)K}\Om_0$.  Every subsequence
has a further subsequence on which $F_1^{(m)}\to F_1$ almost everywhere.  On it,
$|F_1^{(m)}-F_1|^2|y|^2\le4C^2|y|^2\in L^1(\dphi)$, and dominated convergence gives
$(F_1^{(m)}-F_1)y\to0$; hence $(F_1^{(m)}-F_1)y\to0$ along the whole sequence.  The semigroup
is bounded, so $z_m:=e^{-(\tau'-\tau)K}F_1^{(m)}y\to z:=e^{-(\tau'-\tau)K}F_1y$.  Then
$F_2^{(m)}z_m-F_2z=F_2^{(m)}(z_m-z)+(F_2^{(m)}-F_2)z\to0$, by $|F_2^{(m)}|\le C$ and the same
subsequence argument for $z$.  So the left side of (\ref{eq:fknasym}) for the approximants
tends to the left side for $F_1,F_2$.
\item[(4)] \emph{The path side converges.}  $J_\tau$ is an isometry
$L^2(\dphi)\to L^2(\dmu)$ (proof of Lemma~N.22), and $|J_sF_i^{(m)}|\le C$, $|J_sF_i|\le C$ for
$s=\tau,\tau'$ by Step~0.  So
\[
 \bigl\|J_\tau F_1^{(m)}\,J_{\tau'}F_2^{(m)}-J_\tau F_1\,J_{\tau'}F_2\bigr\|_{L^2(\dmu)}
 \le C\bigl(\|F_1^{(m)}-F_1\|_2+\|F_2^{(m)}-F_2\|_2\bigr)\to0 .
\]
Also $e^{-\Uc_{h_{T,g}}}\in L^2(\dmu)$ by A3, and Cauchy--Schwarz gives the convergence of the
right side of (\ref{eq:fknasym}).
\item[(5)] \emph{Conclusion.}  (\ref{eq:fknasym}) holds for $F_1,F_2$.  With $F_1=F_2=1$,
Theorem~V.15 itself gives $\ip{\Om_0}{e^{-2TK}\Om_0}=Z(h_{T,g})$.
\end{enumerate}
So (\ref{eq:fknasym}) rests on Theorem~V.15 as printed, on A1(c), A3 and Lemma~N.22, and on A2's
dictionary; Simon's Remark is not used.  Insert $K=H+E(g)$; the factors $e^{-2TE(g)}$ cancel in
the quotient.  The operators $F_2$ (multiplication by a bounded real function) and
$e^{-(T-\tau')H}$ are self-adjoint, so
\begin{equation}
 \int J_\tau F_1\,J_{\tau'}F_2\,dq_{h_{T,g}}
 =\frac{\ip{F_2\psi_{T-\tau'}}{e^{-(\tau'-\tau)H}F_1\psi_{T+\tau}}}{\|\psi_T\|^2}
 =\frac{\ip{F_1\psi_{T+\tau}}{e^{-(\tau'-\tau)H}F_2\psi_{T-\tau'}}}{\|\psi_T\|^2} ;
 \label{eq:fknbdd}
\end{equation}
the second equality holds because the left side is real, so the middle term equals its complex
conjugate.

\emph{Step 2: truncation.}  Fix a real-valued representative of $X$, and for $\ell>0$ let
$X^{[\ell]}:=(-\ell)\vee(X\wedge\ell)$, the clipped truncation of Lemma~4.3(ii).  It is a
bounded real measurable function of the time-zero field, so (\ref{eq:fknbdd}) holds with
$F_1=F_2=X^{[\ell]}$, and $J_\tau(X^{[\ell]})=(J_\tau X)^{[\ell]}$ by Step~0.  Now let
$\ell\to\infty$.
\begin{itemize}
\item \emph{Path side.}  $(J_\tau X)^{[\ell]}(J_{\tau'}X)^{[\ell]}\to J_\tau X\cdot J_{\tau'}X$
pointwise, with the majorant $|J_\tau X|\,|J_{\tau'}X|$.  The majorant is
$dq_{h_{T,g}}$-integrable.  $X\in L^4(\dphi)$, being a finite sum of Wick powers (bound (V.6b)
quoted in Assumption~A1(d)), and $J_\tau$ is an isometry of $L^4$, so $J_\tau X\in L^4(\dmu)$.
Lemma~N.32 (\texttt{aux:d2-int}; hypothesis $h_{T,g}\in\CE$, Lemma~3.1(ii)) gives
$J_\tau X\in L^2(dq_{h_{T,g}})$, and Cauchy--Schwarz in $L^2(dq_{h_{T,g}})$ gives the
integrability.  By dominated convergence the left side of (\ref{eq:fknbdd}) tends to
$\int J_\tau XJ_{\tau'}X\,dq_{h_{T,g}}$.
\item \emph{Hilbert side.}  For $b>0$, $\psi_b=e^{bE(g)}e^{-bK(g)}\Om_0\in\Ran e^{-bK(g)}\subseteq
D(X)$ by A5(i); so $|X|^2|\psi_b|^2\in L^1(\dphi)$.  Since $|X^{[\ell]}-X|^2|\psi_b|^2\le
|X|^2|\psi_b|^2$ and $X^{[\ell]}\to X$ pointwise, dominated convergence gives
$X^{[\ell]}\psi_b\to X\psi_b$ in $\HH$; this is the computation of Lemma~4.3(ii).  Take
$b=T+\tau$ and $b=T-\tau'$, both $\ge\frac12$; since $e^{-(\tau'-\tau)H}$ is bounded, the
right side of (\ref{eq:fknbdd}) tends to
$\ip{X\psi_{T+\tau}}{e^{-(\tau'-\tau)H}X\psi_{T-\tau'}}/\|\psi_T\|^2$.
\end{itemize}
This proves (ii).  Simon's Remark on p.~164 (``More generally, the above holds if
$G_j^F\in\bigcap_{p<\infty}L^p(Q_F)$ and $G_j^{s_j}=J_{s_j}G_j$'') would give (ii) directly; the
proof above uses it nowhere, neither for $V(h)$ nor for the bounded truncations.

(iii) By Lemma~\ref{lem:Vbar}(i) with $h'=h$, $\Vb(h)=\int_{-1/2}^{1/2}J_\tau X\,d\tau$ in
$L^2(\dmu)$.  For $W\in L^4(\dmu)$ the map $U\mapsto\int UW\,dq_{h_{T,g}}$ is a bounded linear
functional on $L^2(\dmu)$: $|\int UWdq_{h_{T,g}}|\le
Z(h_{T,g})^{-1}\|U\|_2\|W\|_4\|e^{-\Uc_{h_{T,g}}}\|_4$, finite because $e^{-4\Uc_h}\in L^1(\dmu)$
(Assumption~A3).  $\Vb(h)$ and every $J_\tau X$ lie in $L^4(\dmu)$ (Lemma~\ref{lem:Vbar}(ii);
$J_\tau$ is isometric on $L^4$).  Applying the functional with $W=\Vb(h)$, and then for each
$\tau$ with $W=J_\tau X$,
\[
 \int\Vb(h)^2dq_{h_{T,g}}=\int_{-1/2}^{1/2}d\tau\int_{-1/2}^{1/2}d\tau'\,
 \int J_\tau XJ_{\tau'}X\,dq_{h_{T,g}} .
\]
\emph{Dominator.}  For $b\ge\frac12$, $X\psi_b=Xe^{-H/2}\psi_{b-1/2}$ and $\|\psi_{b-1/2}\|\le1$,
so $\|X\psi_b\|\le\|Xe^{-H/2}\|<\infty$ by (i).  In (ii), $T+\tau\ge\frac12$ and
$T-\tau'\ge\frac12$, and $\|\psi_T\|\ge c_g$ by Lemma~4.4 (\texttt{lem:slabvac}); so the
integrand is bounded by $\|Xe^{-H/2}\|^2c_g^{-2}$, uniformly in $T\ge1$ and $\tau\ne\tau'$
(by (ii) and the symmetry of the integrand).  \emph{Limit.}  Lemma~4.4 gives
$\psi_s=\|\psi_s\|\hat\psi_s\to c_g\Om$; hence $X\psi_{T+\tau}=Xe^{-H/2}\psi_{T+\tau-1/2}\to
c_gXe^{-H/2}\Om=c_gX\Om$ ($H\Om=0$), likewise $X\psi_{T-\tau'}\to c_gX\Om$, and
$\|\psi_T\|^2\to c_g^2$; so the integrand tends to $\ip{X\Om}{e^{-|\tau-\tau'|H}X\Om}$ for
$\tau\ne\tau'$.  The integrand is continuous off the diagonal (by (ii): $b\mapsto X\psi_b$ is
continuous on $[\frac12,\infty)$ by (i), and $e^{-uH}$ is strongly continuous), hence
measurable, and the diagonal is a null set.  Dominated convergence and Tonelli (with the spectral
measure $\mu_{X\Om}$ of $X\Om$, and $e^{-|\tau-\tau'|E}\ge0$) give
\[
 \lim_{T\to\infty}\int\Vb(h)^2dq_{h_{T,g}}=\iint\ip{X\Om}{e^{-|\tau-\tau'|H}X\Om}\,d\tau\,d\tau'
 =\int k(E)\,d\mu_{X\Om}(E)=\ip{X\Om}{k(H)X\Om} .\qedhere
\]
\end{proof}

\emph{The role of the cutoff $h_{T,g}$} (item (c) of the commission).  Theorem~\ref{thm:TD4}
holds for every $h\in\CE$; it is used only at the slab cutoffs $h_{T,g}$, $T\ge1$, $g\in\CH$,
which lie in $\CE$ by Lemma~3.1(ii).  The limit $T\to\infty$ is taken at fixed $g$, with the
$g$-dependent dominator $\|V(h)e^{-H(g)/2}\|^2c_g^{-2}$; that is harmless because the bound
$M_B^2$ does not depend on $T$ or $g$.  The condition $T\ge1$ keeps the window
$[-\frac12,\frac12]$ at distance at least $\frac12$ from the ends of the slab, which is what the
dominator needs.  The slab-transfer statement that carries the passage is A2's multi-time
formula for bounded functions at two times placed asymmetrically in the slab (Step~1 of the
proof of Lemma~\ref{lem:fkn}(ii)).  It is extended to $V(h)$ by the truncation of the record's
Lemma~4.3(ii).  The record's Lemma~4.1 states the two-time formula only at the symmetric times
$\mp t/2$, and its one-time formula (3) at any time $s$.

\begin{ttheorem}[bound on the resolvent-weighted moment]\label{thm:Mt}
For every $\lambda_0\ge1$, $\Mt(\lambda_0)\le c_k^{1/2}M_B$.
\end{ttheorem}

\begin{proof}
Fix $g\in\CH$ and $h$ as in Definition~\ref{def:Mt}, choose $a$ with $\supp h\subseteq[a,a+2]$,
and let $X=V(h)$, $H=H(g)$, $\Om=\Om_g$, $Q=Q_g$.  Let $\mu$ be the spectral measure of $QX\Om$
for $H$; it is supported in $\spec H\subseteq[0,\infty)$.  By Lemma~\ref{lem:ineq},
\begin{align*}
 \|(\lambda_0-iH)^{-1}QX\Om\|^2=\int\frac{d\mu(E)}{\lambda_0^2+E^2}&\le c_k\int k\,d\mu\\
 &=c_k\ip{QX\Om}{k(H)QX\Om}\le c_k\ip{X\Om}{k(H)X\Om},
\end{align*}
the last step because $P_0$ commutes with $k(H)$ and $k\ge0$:
$\ip{X\Om}{k(H)X\Om}=\ip{QX\Om}{k(H)QX\Om}+\ip{P_0X\Om}{k(H)P_0X\Om}$.  By
Lemma~\ref{lem:fkn}(iii) and Theorem~\ref{thm:TD4} at $h_{T,g}\in\CE$ with $h'=h$,
$\ip{X\Om}{k(H)X\Om}=\lim_T\int\Vb(h)^2dq_{h_{T,g}}\le M_B^2$.  Take the supremum over $g$ and
$h$.  No gap is used in this proof.
\end{proof}

\emph{Comparison with the sketch.}  Step (c) states $\Mt^2\le2\sup_T\int Z^2dq_{h_{T,g}}$ for
$Z=\int_{-1/2}^{1/2}J_\tau V(h_j)\,d\tau$ (here $\Vb(h_j)$); this holds, since $c_k\le2$ and the
limit in Lemma~\ref{lem:fkn}(iii) is at most the supremum.  Step (d) states that the proof of
Theorem~D.4 applies to $Z^2$ ``with the same box $B_{3\times5}$, the same rings and the same
constant $M_B^2$, since $\|Z\|_2\le\|V(h')\|_2$ and the GJS norm is unchanged''.  That is
confirmed by Theorem~\ref{thm:TD4}, with one precision: the GJS norm of the decomposition
changes, its bound does not.

\section{Assembly: proof of Theorem~\ref{thm:routeT}}\label{sec:assembly}

(a) is Theorem~\ref{thm:Mt}.

(b) Theorem~7.8 gives (G$_\gamma$) with $\gamma=m_1$; by (a), $M'=\Mt_*$ satisfies
$M'\ge\Mt(\lambda_0)$, so Lemma~\ref{lem:rate} and Proposition~\ref{prop:onepath} hold for
$\Psit_{m_1}$.  \emph{Theorem~11.1.}  Step~1 of its proof: given $\varepsilon>0$, choose $d\ge1$
with $\Psit_{m_1}(d)<\varepsilon$ (Lemma~\ref{lem:rate}), and $g_0\in\G$ whose plateau interior
contains $O_d$; for $g,g'\succcurlyeq g_0$, $g\equiv g'\equiv1$ on $O_d$, and
Proposition~\ref{prop:onepath} gives $\|(\omega_g-\omega_{g'})|_{\A(O)}\|<\varepsilon$.  Step~2
(limit functionals, compatibility, the state $\omega_\infty$) uses nothing else and is
unchanged.  Step~3 (the rate): for $g\equiv1$ on $O_d$ and every $g'$ in the cofinal set
$\{g'\in\G:g'\equiv1\text{ on }O_d\}$, Lemma~\ref{lem:rate} gives
$\|(\omega_g-\omega_{g'})|_{\A(O)}\|\le\Psit_{m_1}(d)$, and the bound passes to the norm limit.
The limit state is the one of the record: norm limits are unique.  \emph{Proposition~11.2} uses
only the convergence of Theorem~11.1.  \emph{Theorem~12.2:} the record's proof derives (1)--(2)
from Theorem~11.1 and Proposition~11.2, (3) from (1) and Proposition~12.1, and (4) from (1),
Lemma~2.8 and Proposition~E.5; with the modified Theorem~11.1 these derivations are unchanged.
\emph{Corollary~12.3:} its proof uses (1) in the form $\Psi_{m_1}(2n-R)\to0$, which holds for
$\Psit_{m_1}$.

(c) The statements used are those listed in (a)--(b), Lemmas~\ref{lem:transforms}--\ref{lem:rate}
and~\ref{lem:ineq}--\ref{lem:fkn}, Theorems~\ref{thm:TD4}--\ref{thm:Mt}, and, from the record,
Theorem~7.8, Propositions~2.6, 2.7, 8.1, 8.2, Lemma~9.1, Step~2 of Lemma~9.3, the first display
of the proof of Lemma~10.1, the first assertion of Lemma~10.2, Step~2 of Theorem~11.1,
Propositions~11.2, 12.1, Lemma~2.8, Proposition~E.5, the inputs of Theorem~\ref{thm:TD4}, and
the inputs of Lemma~\ref{lem:fkn}: Lemmas~3.1(ii), 4.4, N.22 and N.32, the computation of
Lemma~4.3(ii), and Assumptions A1(c), A1(d), A2 (with Simon's Theorem~V.15 as printed behind it),
A3 and A5(i).  In the record, Theorem~5.1 enters the chain in two places.
\begin{itemize}
\item Through the second assertion of Lemma~10.2 (l.~2776), hence Lemma~10.3,
Proposition~10.4, Lemma~10.5 and the instantiation of $\Psi_{m_1}$ at the head of \S11
(l.~2921).
\item Through the first sentences of the proof of Theorem~12.2 (ll.~3104--3106): ``(5) is
Theorem~7.8 and Theorem~5.1.  These supply (G$_\gamma$) with $\gamma=m_1$ and (M$_2$), so
Theorem~11.1 and Proposition~11.2 give (1)--(2)''.  So (M$_2$) reaches items (1)--(2) through
Theorem~11.1, and item (5) directly.
\end{itemize}
None of these is used for items (1)--(4) here: Theorem~11.1 and Proposition~11.2 are re-proved
with Proposition~\ref{prop:onepath} and Lemma~\ref{lem:rate}, and only the moment clause of
item (5) still takes Theorem~5.1.  Theorem~D.4 enters only through Step~5 of Theorem~5.1, and
A8(iv)-T only through Theorem~D.4 (Assumption~A8, item (iv-T): ``Used in Theorem~D.4 (through
Lemma~D.2), hence in Theorem~5.1 and in every statement that consumes its constant $M$ \dots\
not in \S7'').  A search of \S7 of the record (ll.~2055--2389) for \texttt{thm:M2},
\texttt{thm:D4}, \texttt{M\_B} and ``transverse'' finds nothing.  The applications of
Theorem~1.1.7 are those of Theorem~7.8 (through Lemma~7.4 and Theorem~D.3; cases P1, P2) and of
Theorem~\ref{thm:TD4} (cases P3, P4).  \qed

\section{The clustering input actually used}\label{sec:input}

\subsection{What GJS print}\label{sec:printed}

The following sentences were read on the rendered pages (images named in \S\ref{sec:setting});
nothing longer is needed.
\begin{itemize}
\item p.~590: ``\textsc{Theorem} 1.1.1. \emph{Let $\lambda/m_0^2$ be sufficiently small.  The
$\mathscr P(\phi)_2$ quantum field model exists as a limit of space cutoff theories and satisfies
the Wightman axioms.}'' (the statement continues).
\item p.~593: ``(1.1.5) $Q(x_0)=\int d\vec x\,w(x)\prod_{i=1}^r{:}\Phi(x_i)^{n_i}{:}$'' and
``(1.1.6) $Q=\int dx_0Q(x_0)$''; ``$Q$ is said to be localized if $w(x)$ is supported in a
product $\Delta_{j_1}\times\cdots\times\Delta_{j_n}\subset R^{2r}$ where $\Delta_{j_i}$ is the unit
lattice square in $R^2$ centered at the lattice side $j_i\in Z^2$.''
\item p.~594: ``For two localized monomials $Q_1$ and $Q_2$, with disjoint localization squares
$\{\Delta^{(1)}_{j_i}\}$ and $\{\Delta^{(2)}_{j_l}\}$ respectively, let $d$ be the width of the
widest strip bounded by two parallel lines, and separating the $\{\Delta^{(1)}_{j_i}\}$ from the
$\{\Delta^{(2)}_{j_l}\}$'' (no period in the source).  ``\textsc{Theorem} 1.1.7. \emph{Let $m_0$,
$K_2$ and $K_1$ be sufficiently large and let $\lambda$ be sufficiently small.  There are constants
$0<m_1$ and $O(1)$ independent of $h$, $Q_1$ and $Q_2$ such that} (1.1.10)
$|\int Q_1Q_2dq_h-\int Q_1dq_h\int Q_2dq_h|\le O(1)e^{-m_1d}\|Q_1\|\,\|Q_2\|$.''
\item p.~595, in the proof of (1.1.11) with ``$Q=\prod_i\Phi(f_i)$'' for ``test functions
$f_i\in C_0^\infty(\mathcal O)$'' (p.~594) and $V(g_i)$, $g_i=g\chi_{\Delta_i}$: ``We apply
Theorem 1.1.7 to each term in the sum to obtain'' a sum over ``$i\in Z^2\cap\operatorname{supp}g$''
of ``$O(1)e^{-m_1\operatorname{dist}(\mathcal O,\Delta_i)}$'', followed by ``As a special case,
the cutoff Schwinger functions $S_h$ converge as $h\to1$.''  Then: ``In the monomial $Q(x_0)$
of (1.1.5), suppose $x_0=(x_{1,0},\cdots,x_{r,0})$ where $x_{1,0}\le\cdots\le x_{r,0}$.''  And
``\textsc{Theorem} 1.1.8. \emph{For $Q(x_0)$ of the form} (1.1.5), $|\int Q(x_0)dq_h|\le
O(1)\|Q(x_0)\|$, \emph{uniformly as $h\to1$}.''
\item p.~596: ``We let $A$ and $B$ denote monomials of type (1.1.5) of degree $|A|$ and $|B|$ and
with kernels $w_A$ and $w_B$ respectively.  We assume that $A$ is localized at times
$x_{i0}\le0$, and that $B$ is localized at times $x_{i0}\ge0$.''  ``In the next two results, we
take all times $x_{1,0}=t=\cdots=x_{r,0}$ equal in $Q(x_0)$.  To indicate this restriction, we
write $Q(t)$ in place of $Q(x_0)$.''  Theorem~1.1.10 then names an approximant
``$Q(T)=\alpha(T)+\int\Phi(0,\vec x)h_T(\vec x)d\vec x$'', and Theorem~1.1.11 ``$n$\textsuperscript{th}
degree polynomials $Q_n(t_j)$ of the form (1.1.5) with $r=1$'' with the bound ``(1.1.16)
$|\int(A-Q)B_Tdq_h|\le O(1)\|A\|\,\|B\|e^{-\gamma T}$, uniformly in $h$''.  The proof of
Corollary~1.1.9 on the same page writes ``Apply the theorem to
$Q_1=Q(x_{i0}-x_{10})Q(x_{10}-x_{i0})^-$, where $Q$ has the form (1.1.5)''; that ``the theorem''
is Theorem~1.1.8 is a reading, supported by the form of the bound that follows,
$\int Q_1dq\le O(1)\|Q_1\|_{K_2}$.
\item p.~597: ``After rotating and translating the strip separating $Q_1$ and $Q_2$ in Theorem
1.1.7, and after choosing $A=Q_1$, $B_T=Q_2$ in Theorem 1.1.11, the proof follows from Theorem
1.1.8, (1.1.16) and the above determination of $Q$.''  In the proof of Theorem~1.1.5 on the same
page: ``where $\theta_A$ is the vector (1.1.13) corresponding to $A$''.
\item p.~600: ``Proof of Theorem 1.1.8, assuming Theorem 1.1.10.  Choose $T=0$ and $B=I$ in
Theorem 1.1.11.  Then $Q=0$, so Theorem 1.1.8 is a special case of Theorem 1.1.10.''  Both
occurrences of ``1.1.10'' are misprints for 1.1.11; the record marks them [sic] (A8(vi),
ll.~1236--1238).
\item p.~628, ``Proof of Theorems 1.1.8 and 1.1.11'': the expansion (4.3) groups ``terms with
common values of $i$ and $y_0$'', and ``set $Q_n(y_0)=\sum_{i=0}^n\int d\vec y\,w_i(y)\int
{:}\Phi(y_1)\cdots\Phi(y_i){:}$\,.  (By removal of the Wick ordering in $Q_n(y_0)$, we may apply
the norms of (1.1.8).)''
\item p.~629: ``Theorem 1.1.11 follows from this inequality, while Theorem 1.1.8 is the special
case $B=I$, $n=0$.''
\end{itemize}

\subsection{The applications of Theorem 1.1.7 behind items (1)--(4)}\label{sec:cases}

Before Route T the record applies Theorem~1.1.7, for items (1)--(4), in the cases P1, P2 and in
Theorem~D.4; after Route T in the cases P1--P4.  All cutoffs are in $\CE$, and the constant used
is the printed $O(1)=C_7$ ``independent of $h$'', read over $\CE$ (A8 Scope (d)).

\begin{center}\small
\begin{tabular}{@{}p{0.05\textwidth}p{0.15\textwidth}p{0.22\textwidth}p{0.17\textwidth}p{0.12\textwidth}p{0.13\textwidth}@{}}
\toprule
case & where & $Q_1$ & $Q_2$ & strip & cutoffs\\
\midrule
P1 & Lemma 7.4 & sharp-time cylinder pieces at $-t/2$, type (1.1.5), one time & the same at
$t/2$ & parallel & every $h\in\CE$\\
P2 & Theorem D.3 (Lemma D.2) & sharp-time cylinder pieces at $s$, type (1.1.5), one time &
$U_{\Delta,n}$, type (1.1.6), $r=1$ & parallel ($P_l$) & $h_\alpha$, $h=h_{T,g}$\\
P3 & Theorem~\ref{thm:TD4}, class (P) & $\overline Q_{n,n',k,k'}$, type (1.1.6), $r=2$ &
$U_{\Delta,n''}$, type (1.1.6), $r=1$ & parallel & $h_\alpha$, $h\in\CE$\\
P4 & Theorem~\ref{thm:TD4}, class (T) & as in P3 & as in P3 & lines of constant $x^1$ & as in
P3\\
\midrule
--- & Theorem D.4, class (T) (A8(iv)-T; not used after Route T) & $Q_{n,n',k,k'}$, type (1.1.5),
two factors at time $0$ & $U_{\Delta,n''}$ & lines of constant $x^1$ & $h_\alpha$, $h\in\CE$\\
\bottomrule
\end{tabular}
\end{center}

Theorem~D.4's class (P) applications, of the form P2 with two factors at time $0$, are not used
after Route T either.

\subsection{Two readings of p.~596}\label{sec:readings}

The sentence ``In the next two results, we take all times $x_{1,0}=t=\cdots=x_{r,0}$ equal in
$Q(x_0)$'' admits two readings.
\begin{itemize}
\item[\textbf{R1}] It fixes the form of the approximants: the notation $Q(t)$ it introduces
occurs in Theorems~1.1.10--1.1.11 only for $Q(T)$ and $Q_n(t_j)$.  The monomials $A$ and $B$
are monomials of type (1.1.5) in the sense of p.~595, whose factors may sit at different times
$x_{1,0}\le\cdots\le x_{r,0}$.
\item[\textbf{R2}] It restricts $A$ and $B$ themselves to monomials whose factors all sit at one
time.
\end{itemize}
\emph{In favour of R1} (printed passages of \S\ref{sec:printed}):
\begin{enumerate}
\item[(a)] p.~595 introduces the times $x_{1,0}\le\cdots\le x_{r,0}$ of a monomial of type
(1.1.5) and states Theorem~1.1.8 ``For $Q(x_0)$ of the form (1.1.5)''.  The proof of
Corollary~1.1.9 applies ``the theorem'' (Theorem~1.1.8, by the form of the bound) to a product
whose factors sit at different times.  On p.~600 GJS prove Theorem~1.1.8 by choosing $T=0$ and
$B=I$ in Theorem~1.1.11, and p.~629 repeats that ``Theorem 1.1.8 is the special case $B=I$,
$n=0$''.  So GJS obtain Theorem~1.1.8 as the case $A=Q(x_0)$, $B=I$ of Theorem~1.1.11.  On R2
this derivation gives Theorem~1.1.8 only for monomials with all their factors at one time, short
of the scope at which it is printed and used; on R1 it gives the printed scope.
\item[(b)] p.~596 says ``$A$ is localized at times $x_{i0}\le0$'', with an index on the times.
\item[(c)] In the proof of Theorem~1.1.5 on p.~597, ``$\theta_A$ is the vector (1.1.13)
corresponding to $A$''; (1.1.13) is the vector that p.~595 builds for ordered, possibly distinct
times (weaker evidence).
\item[(d)] GJS's own application of Theorem~1.1.7 on p.~595 needs R1 (below).
\end{enumerate}
\emph{In favour of R2:} only the wording ``in $Q(x_0)$'', which could be read as concerning every
monomial of the next two results.  The draft of this note offered a second argument, that the
approximants have one Wick factor ($r=1$ in the statement of Theorem~1.1.11), so the sentence
would be redundant for them.  p.~628 removes its premise: the approximants are built as
$Q_n(y_0)$, with up to $n$ field factors at one time $y_0$ after the removal of the Wick
ordering.  So the sentence has content for them.  No printed sentence states in words that $A$
may have factors at different times, and in that narrow sense the printed text does not decide
between R1 and R2.  But R1 is the reading under which GJS's own proofs are complete: at
p.~595, and in the derivation of Theorem~1.1.8 at its printed scope (pp.~600, 629).

\emph{Provenance.}  R1 is the ``multi-time reading of (1.1.5)'' that Remark~N.39 of the record
names (l.~6816), applied to the monomial $A$ of Theorem~1.1.11.  The record's own objects
already carry it: Definition~N.26(b) admits times ``equal or not'' (l.~5750), and Remark~N.27
says ``Case (b) admits distinct times $x_{i,0}$, the reading of p.~595'' (l.~5783).
\path{ugjs/PLAN.md}, Part~2 (ll.~312--315 at 49f2d23), records that monomials of type (1.1.6)
``are reached in any orientation, after slicing, under the multi-time reading of (1.1.5)'',
and that ``This is how GJS's own use on p.~595 is covered''.

\emph{GJS's own use needs R1.}  On p.~595 GJS apply Theorem~1.1.7 to the pairs
$(\prod_i\Phi(f_i),V(g_i))$ for all $i\in Z^2\cap\supp g$, that is, across strips of every
orientation relative to $\mathcal O$.  $\prod_i\Phi(f_i)$ is a time-integrated monomial with
$r$ factors of degree $1$ (type (1.1.6), split into localized pieces by (1.1.9)), and $V(g_i)$ is
a sum of time-integrated monomials with $r=1$.  The printed proof takes ``$A=Q_1$'' of type
(1.1.5); for a time-integrated $Q_1$ this is possible only after slicing in time, and for $r\ge2$
the slices have their factors at different times.  So on R2 the printed proof of Theorem~1.1.7
does not reach the application GJS make on p.~595, in the proof of their Theorem~1.1.1 (p.~590),
at the step where the cutoff Schwinger functions converge (p.~595).  That step needs products
of at least two fields.  On R1 the printed proof reaches it.  This is an observation about the
source's logic; neither the record nor this note verifies any proof of GJS.

\subsection{What the printed proof reaches}

Following the usage of the record (A8 Scope (b), Remark~N.39, \S D.e.4), a pair
$(Q_1,Q_2)$ separated by a strip $S$ is said to be \emph{reached by the printed proof} of
Theorem~1.1.7 if, after the rigid motion that the proof performs (``rotating and translating
the strip''), one of the following holds.
\begin{itemize}
\item[(a)] $Q_1$ and $Q_2$ are admissible as $A$ and $B_T$ in Theorem~1.1.11, and $B_T$ as
$Q(x_0)$ in Theorem~1.1.8.  Relabelling $Q_1\leftrightarrow Q_2$ is allowed, since (1.1.10) is
symmetric in $Q_1$ and $Q_2$.
\item[(b)] $Q_1$ and $Q_2$ are integrals over time of slices for which (a) holds pairwise, and
the bound for $(Q_1,Q_2)$ follows from the bounds for the pairs of slices by an integration that
is written out (as in Proposition~\ref{prop:N39}).
\end{itemize}
The notion concerns the scope of the printed proof, not its correctness, which neither the record
nor this note examines (A8 Scope (c)).  Three questions concern the motions and the constant of
the printed proof.
\begin{itemize}
\item \emph{The constant.}  The printed proof takes its constant ``from Theorem 1.1.8'' at the
same $h$ (p.~597).  p.~595 prints Theorem~1.1.8 ``uniformly as $h\to1$'', while p.~629 states that
``the estimates of Theorems 1.1.8, 1.1.11, 3.1 and 4.1 are uniform in the space cutoff $h$''
(Source Fact~2.4, \texttt{sf:printed}).  This question is common to P1--P4 and to every use of
Theorem~1.1.7 in the record; it is not decided here.
\item \emph{Translations} raise no question in P1--P3; in P4 they raise at most a bounded
factor, of the same kind as the question of the constant.  In P1 and P2 both partners are
sharp-time pieces.  In P3 and P4 the printed proof is applied after slicing, so again to
sharp-time pieces: one time per slice of $U_{\Delta,n}$, two times per slice of a two-factor
piece.
  \begin{itemize}
  \item A translation by $a\in\R^2$ is covariant.  Lemma~N.13: $T_a$ preserves $d\mu_0$, $\tau_a$
  maps $\CE$ onto $\CE$, and $\int F\circ T_a\,dq_{\tau_ah}=\int F\,dq_h$.  It also carries a
  constructed monomial to the constructed monomial with the translated kernel.  Indeed
  $\Phi_\kappa(z)\circ T_a=\Phi_\kappa(z+a)$ (proof of Lemma~N.13(iii)) and the variance
  $c_\kappa$ does not depend on $z$, so the regularized monomial of Definition~N.26 with kernel
  $w$, composed with $T_a$, is pointwise the one with kernel $w(\cdot-a)$.  The limits agree,
  because $F\mapsto F\circ T_a$ is an isometry of $L^2(\dmu)$ (Lemma~N.13(i)).
  \item A sharp-time factor moved to another time lies in the square of the row of its new time.
  Remark~N.2(ii): ``The choice of row changes neither the local degrees nor
  $\|w(x_0,\cdot)\|_2$, hence not the norm (1.1.8)''.
  \item In P3, both factors of a two-time slice start in row $0$.  A translation can therefore
  only split them into different rows of one column.  That splits a local degree $n+n'$ into
  $n$ and $n'$, which does not increase the norm (1.1.8): $K_2^{2(n+n')}$ is unchanged and
  $n!\,n'!\le(n+n')!$.
  \item In P4, after the rotation, the two factors of a piece with $k\ne k'$ start in different
  rows of one column, and over part of the range their times are less than $1$ apart.  A
  non-integer translation can then put both into one square.  That raises the weight in (1.1.8)
  by $(n+n')!/(n!\,n'!)\le2^{n+n'}\le2^{4p'}$.  Translations by integers avoid this: they map
  the lattice onto itself and change no norm, at the price of at most one unit of the separation
  $T$ in (1.1.16), a factor at most $e^{\gamma}$.  Either way the effect is a bounded factor, of
  the same kind as the question of the constant; the verdicts below are not affected.
  \end{itemize}
\item \emph{Rotations} are needed only for P4 (and for the record's class (T) of
Theorem~D.4).  For P1--P3 the strip is parallel, and relabelling $Q_1\leftrightarrow Q_2$ replaces
the rotation by $\pi$.  A rotation $\varrho$ by $\pi/2$ about a lattice point maps the squares of
p.~593 onto themselves, and lines of constant $x^1$ onto lines of constant $x^0$.  ``Rotating''
presupposes that it also carries $dq_h$ to $dq_{h\circ\varrho^{-1}}$, and the constructed
monomials to those with the rotated kernels.  Call this the \emph{rotation presupposition}.
  \begin{itemize}
  \item The record proves translation covariance (Lemma~N.13), not rotation covariance.
  \item A proof would follow Lemma~N.13 with a radial mollifier.  It would use that the kernel $G$
  of the free covariance is radial (Lemma~N.19(v)) and that the limits do not depend on the
  mollifier (Lemma~N.28(e)).
  \item It is not written here, and no proof in the chain of Theorem~\ref{thm:routeT} uses it.
  \end{itemize}
\end{itemize}

\begin{tprop}[what the printed proof reaches]\label{prop:reach}
\begin{enumerate}
\item[\rm(a)] The pairs of case P1 are reached, on R1 and on R2.
\item[\rm(b)] The pairs of case P2, and those of Theorem~D.4's class (P), are reached after the
time slicing of $U_{\Delta,n}$, on R1 and on R2 (Proposition~\ref{prop:N39}).
\item[\rm(c)] The pairs of case P3 are reached on R1.  The pairs of case P4 are reached on R1,
granted the rotation presupposition.  On R2, the pairs of P3 and P4 whose two-factor piece
$\overline Q_{n,n',k,k'}$ has $n,n'\ge1$ are not reached.  The pieces with a factor of degree $0$
are reached after slicing, on either reading (for P4, granted the rotation presupposition).
\item[\rm(d)] The pairs of Theorem~D.4's class (T) (A8(iv)-T) are reached neither on R1 nor
on R2.
\end{enumerate}
\end{tprop}

\emph{Status.}  This is a statement about the scope of the source's printed proof, with the notion
of ``reached'' just defined.  Proved: (a), (b), (d); in (c), the assertions for P3 on R1, and
the assertions on R2 except the reach of P4's degree-$0$ pieces.  Proved conditional on the
rotation presupposition: in (c), P4 on R1, and the reach of P4's degree-$0$ pieces.

\begin{proof}
(a) Both monomials are of type (1.1.5), each with all its factors at one time, and the strip is
parallel.  After a time translation, and relabelling $Q_1\leftrightarrow Q_2$ if necessary, the
earlier monomial is localized at times $\le0$ and the later at times $\ge$ the width of the strip.
On either reading they are admissible as $A$ and $B_T$, and $B_T$ in Theorem~1.1.8.

(b) Proposition~\ref{prop:N39}.

(c) \emph{P3 on R1.}  Fix the following data:
\begin{itemize}
\item a piece $\overline Q:=\overline Q_{n,n',k,k'}$ of Lemma~\ref{lem:Vbar}(iii);
\item a square $\Delta=\Delta_{l,k''}\in\mathcal D(h,B)\cap\mathcal L_j$ of class (P), so that
$|l|=1+j$;
\item a degree $1\le n''\le2p'$, and $\alpha\in[0,1]$.
\end{itemize}
Write $U:=U_{\Delta,n''}$ and $q:=dq_{h_\alpha}$.

\emph{Step 1: two-time slices.}
\begin{itemize}
\item By Lemma~\ref{lem:slice}, $\Vb_{n,k}=\int\Vb_{n,k}(\tau)\,d\tau$ in $L^2(\dmu)$.  The slices
$\Vb_{n,k}(\tau)=\lambda a_nJ_\tau W_n(h'_k)$ lie in the $n$-th chaos.  This follows from the
proof of Lemma~\ref{lem:Vbar}(i), since $J_\tau$ maps a Wick power to the Wick power of the image
variable (proof of Lemma~N.22).
\item On a finite sum of chaoses $\|\cdot\|_4\le3^{p'}\|\cdot\|_2$, by the argument of
Lemma~N.21(ii).  So $\tau\mapsto\Vb_{n,k}(\tau)$ is also strongly measurable and Bochner
integrable in $L^4(\dmu)$, with the same integral.
\item The product $L^4\times L^4\to L^2$ is bounded and bilinear (H\"older).  So
\[
 \overline Q=\Vb_{n,k}\Vb_{n',k'}=\iint\overline Q(\tau,\tau')\,d\tau\,d\tau'
\]
as an $L^2(\dmu)$-valued Bochner integral, where
$\overline Q(\tau,\tau'):=\Vb_{n,k}(\tau)\Vb_{n',k'}(\tau')$ lies in $L^4(\dmu)$.
\item For $\tau,\tau'\in[-\frac12,\frac12)$, Lemma~N.28(c) (type (b)) identifies
$\overline Q(\tau,\tau')$ with the localized monomial of type (1.1.5) with kernel
$\lambda^2a_na_{n'}h'_k\otimes h'_{k'}$ and two factors at the times $\tau,\tau'$.  It is localized
in $\Delta_{0,k}$ and $\Delta_{0,k'}$, with the local degrees of $\overline Q$.
\item Hence
$\GJS{\overline Q(\tau,\tau')}=K_1\bigl[\prod_\Delta K_2^{2N(\Delta)}N(\Delta)!\bigr]
\lambda^2|a_na_{n'}|\,\|h'_k\|_2\|h'_{k'}\|_2$.  Since the time factor of the kernel has $L^2$ norm
$1$, $\iint_{[-\frac12,\frac12)^2}\GJS{\overline Q(\tau,\tau')}\,d\tau\,d\tau'=\GJS{\overline Q}$.
\end{itemize}

\emph{Step 2: the covariance.}  For $W\in L^4(\dmu)$ the maps $V\mapsto\Cov_q(V,W)$ and
$V\mapsto\Cov_q(W,V)$ are bounded linear functionals on $L^2(\dmu)$:
\[
 \Bigl|\int VW\,dq\Bigr|\le Z(h_\alpha)^{-1}\|V\|_2\|W\|_4\|e^{-\Uc_{h_\alpha}}\|_4,\qquad
 \Bigl|\int V\,dq\Bigr|\le Z(h_\alpha)^{-1}\|V\|_2\|e^{-\Uc_{h_\alpha}}\|_2 ,
\]
both finite by A3.  Apply them first with $W=U\in L^4(\dmu)$, then, for each $(\tau,\tau')$,
with $W=\overline Q(\tau,\tau')$.  With $U=\int U(t)\,dt$ (Lemma~\ref{lem:slice}) this gives
\[
 \Cov_q(\overline Q,U)=\iint d\tau\,d\tau'\int dt\ \Cov_q\bigl(\overline Q(\tau,\tau'),U(t)\bigr).
\]

\emph{Step 3: one triple.}  Let $\tau\ne\tau'$ and let $t$ lie in row $l$.
\begin{itemize}
\item $\overline Q(\tau,\tau')$ is a localized monomial of type (1.1.5) with its factors at the two
times $\tau,\tau'$ of row $0$.  $U(t)$ is one with a single factor at the time $t$, localized in
$\Delta$.
\item The parallel strip of class (P), of width $j$, separates their localization squares.
\item On R1 the pair is reached in the sense of (a).  After a time translation and, if $l<0$,
relabelling, the earlier monomial is admissible as $A$ and the later as $B_T$ in
Theorem~1.1.11, and $B_T$ in Theorem~1.1.8.
\item Its printed bound is $C_7e^{-m_1j}\GJS{\overline Q(\tau,\tau')}\GJS{U(t)}$.
\end{itemize}

\emph{Step 4: integration.}  Steps~1--3 and Lemma~\ref{lem:slice}(ii) give
\[
 |\Cov_q(\overline Q,U)|\le C_7e^{-m_1j}\GJS{\overline Q}\GJS{U} .
\]
This is the bound that
Theorem~\ref{thm:TD4} takes from Lemma~D.2(a) for this pair.  So P3 is reached on R1, in the sense
of (b).

\emph{P4 on R1, granted the rotation presupposition.}  After a rotation $\varrho$ by $\pi/2$ about
a lattice point, the following hold.
\begin{itemize}
\item The strip of class (T) becomes parallel.
\item $\overline Q$ and $U$ become time-integrated monomials with the rotated kernels.
$w_{n,k}\circ\varrho^{-1}$ is again bounded and vanishes outside one square; the product
structure, the $L^2$ norms and the local degrees are preserved.
\item The measure becomes $dq_{h_\alpha\circ\varrho^{-1}}$, and $h_\alpha\circ\varrho^{-1}\in\CE$.
\end{itemize}
Granted that the constructed objects and the measures transform in this way, Steps~1--4 apply in
the rotated frame; the new time of a factor is, up to sign, its old space coordinate.

\emph{On R2.}
\begin{itemize}
\item If $n,n'\ge1$, R2 admits as $A$ or $B$ only monomials with all their factors at one time.
The slices $\overline Q(\tau,\tau')$ with $\tau\ne\tau'$ are therefore admissible neither as $A$
nor as $B$.
\item They carry the whole integral of Step~1, because the diagonal $\tau=\tau'$ is a null set.
So the printed proof gives no bound for them, and the pair is not reached.
\item For P4 the same holds in the rotated frame.  If the rotation presupposition fails, the
printed proof does not reach P4 at all; so P4 is not reached on R2 either way.
\item Pieces with $n,n'\ge1$ occur, because $a_{2p'}>0$.
\item If $n=0$ (or $n'=0$), the factor of degree $0$ is the constant $\lambda a_0\int h'_k$.  So
$\overline Q_{0,n',k,k'}=(\lambda a_0\int h'_k)\,\Vb_{n',k'}$ is a constant times a
time-integrated monomial with $r=1$.  Its pairs with $U$ are reached after slicing both
partners, on either reading.  Both partners are time-integrated with one factor, so all their
slices are one-time monomials, and Steps~1--4 above apply with one-factor slices in place of
two-time slices.  For P4 this holds in the rotated frame, granted the rotation
presupposition.
\end{itemize}

(d) This is the record's analysis (\S D.e.4, \texttt{sub:d4-transverse}, p.~101): a rotation by a
right angle ``turns the sharp-time piece of $Q_V$ into an object with fixed spatial position and
integrated times, which is not of type (1.1.5) in the rotated frame''.  It is not of type (1.1.6)
either, its kernel being concentrated on a line.  Slicing does not help: the slices in the new time
have spatial kernels concentrated at one point, which are not in $L^2$.  This holds on both
readings.
\end{proof}

\begin{trem}[the obstruction on R2; an argued observation]\label{rem:obstruction}
On R2, a pair across a strip bounded by lines of constant $x^1$ is reached only if, after the
rotation, every slice of each monomial has all its factors at one time.
\begin{itemize}
\item A time-integrated monomial with $r\ge2$ factors of positive degree and a nonzero kernel in
$L^2(\R^{2r})$ never has this property, because the set where two of its time coordinates
coincide is a null set of $\R^{2r}$.
\item Nor does a sharp-time monomial of positive degree and positive spatial extent
(Proposition~\ref{prop:reach}(d)).
\end{itemize}
So on R2 the printed proof reaches a transverse pair only if both monomials are time-integrated,
each with at most one factor of positive degree.

Now take the comparison of Lemma~D.2 around a box $B$ of fixed shape, independent of the cutoff.
Apply it to a quadratic Euclidean observable of $V(h)$, that is, a sum of products of two factors
of positive degree at different points, whose localization squares lie in a row $l'$.
\begin{itemize}
\item The charged squares in the rows $l'-1,l',l'+1$ outside $B$ admit no parallel strip of positive
width.  This is Lemma~N.7(iv), which holds for every rectangle of squares containing the
observable's squares.
\item Such squares occur for suitable $h\in\CE$: for example $h=h_{T,g}$ with $T\ge|l'|+2$ and
$\supp g$ charging a column outside $B$.
\item For $B_{3\times5}$, Remark~N.12(b) shows in addition that uniformity in $h\in\CE$ forces the
transverse strips for the other side squares of the rings.
\end{itemize}
So such a comparison applies Theorem~1.1.7 across transverse strips to two-factor pieces of
positive degrees, and it is not reached on R2.

This is an argued observation, not a theorem: the class of ``routes'' is not defined.  A box that
grows with $\supp h$ escapes the argument, at the price of uniformity in $h$; Theorem~D.3 is of
that kind.  The observation locates the step at which the stronger claim fails for routes of this
type.  That claim is: ``after Route T, items (1)--(4) rest only on cases the printed proof reaches,
whatever the reading''.  Not excluded: a route that avoids quadratic observables, or one that
avoids the comparison method.
\end{trem}

\subsection{Verdict}

\emph{The category.}  The commission distinguishes three categories: a case is (i) reached by the
printed proof, (ii) covered by the printed statement only, or (iii) dependent on a reading.
\begin{itemize}
\item The cases P3--P4 that Route T needs fall in category (iii), dependent on a reading.
\item They are covered by the printed statement of Theorem~1.1.7: two localized monomials, a
strip of any orientation, a constant independent of $h$.
\item They are reached by the printed proof on reading R1 of p.~596, the reading under which
GJS's own proofs are complete (\S\ref{sec:readings}).  For P4 this holds granted the rotation
presupposition.
\item On R2, their pieces of positive degrees are covered by the printed statement only, in both
orientations.
\end{itemize}

\emph{What Route T changes.}
\begin{itemize}
\item \emph{Before.}  Items (1)--(4) rest on A8(iv)-T, which the printed proof reaches on no
reading.  They also rest on P1, P2 and Theorem~D.4's class (P), which it reaches on both
readings.
\item \emph{After.}  They rest on P1 and P2, reached on both readings, and on P3 and P4, reached
on R1 (P4 granted the rotation presupposition).
\item \emph{On R1} every application behind items (1)--(4) is reached.
\item \emph{On R2} neither the record nor Route T achieves that.  The record keeps the unreached
class (T) of Theorem~D.4; Route T has the unreached pieces of positive degrees of P3 (parallel)
and P4 (transverse).
\end{itemize}
Route T therefore does not remove every dependence of items (1)--(4) on unreached cases.  It
replaces a case that no reading reaches by reading R1.  It adds no input to the record's imports
(Lemma~\ref{lem:fkn}).

\emph{Readings and presuppositions that remain for items (1)--(4) after Route T.}
\begin{enumerate}
\item R1 of p.~596, and for P4 the rotation presupposition.  Both matter only for the question
whether the printed proof reaches P3--P4.  Under the record's import standard (the printed
statement at its printed scope) neither is needed.
\item The readings the record declares, with the open items of
\path{part_i/C5_reading_record.md} named there:
  \begin{itemize}
  \item A1 Scope (ll.~797--804): the sources' $V(g)$ is the record's.
  \item A3 Scope (ll.~919--927): Simon's and GJS's interaction is $\Uc_h$.
  \item A6 Scope (ll.~1028--1046): the fields of GJ1 and GJ2.
  \item A8 Scope (d) (ll.~1381--1399): the constructed objects are GJS's monomials, $\Uc_h$ is
  GJS's $V(h)$, ``independent of $h$'' is read over $\CE$, and the $O(1)$ constants depend on the
  printed thresholds only.
  \end{itemize}
Route T adds monomials of type (1.1.6) with $r=2$ to the objects of A8 Scope (d); the record's
Definition~N.26(a) already covers them.
\item The question of the constant, taken from Theorem~1.1.8 at the same $h$.  It affects every
case alike, as in the record.  The translation question named in the draft of this note does not
arise in P1--P3.  In P4 it reduces to a bounded factor of the same kind ($\le2^{4p'}$, or
$\le e^\gamma$ with integer translations; see the ``Translations'' item before
Proposition~\ref{prop:reach}).
\end{enumerate}

\section{Remark N.39 of the record: the parallel pairs are reached}\label{sec:N39}

Remark~N.39 (\texttt{aux:d3-orient}, p.~96, ll.~6804--6825) says of the applications of
Theorem~1.1.7 in Theorem~D.3: ``The partner $U_{\Delta,n}$ is of type (1.1.6); the printed proof
reaches the pair only after $U_{\Delta,n}$ is sliced in time and the bound is integrated over the
slices, the multi-time reading of (1.1.5) that the plan records (Part~2) and GJS do not print.''
It leaves open whether the printed proof reaches the pair.  The finding N3 of 22 November states
that it does, by a Fubini step and with no multi-time reading.  The argument, written out:

\begin{tprop}[the pairs of Theorem D.3 are reached after slicing]\label{prop:N39}
Let the data be those of Step~3 of the proof of Theorem~D.3 (\texttt{thm:D3}, p.~95): $g\in\CH$,
$T>0$, $h=h_{T,g}$, $B=B_s(g,I)$, $R=R_s(I)$; $Q_1:=c_\beta Q_{\beta,\kappa}(s)$, a localized
monomial of type (1.1.5) with $r_\beta$ factors of degree $1$, all at the time $s$, localized in
squares of $R_s(I)$ (Lemma~N.36(ii),(v)); $\Delta=\Delta_{l,k}\in\mathcal D(h,B)$, so that
$|l-l_s|\ge2$ (Lemma~N.38(i)); $1\le n\le2p'$; $\alpha\in[0,1]$ and $h_\alpha=h^B+\alpha(h-h^B)$;
$w_\Delta:=|l-l_s|-1$.  Let $U_{\Delta,n}(t)$ be the slices of $U_{\Delta,n}$ given by
Lemma~\ref{lem:slice} (kernel $\lambda a_n\tilde h\mathbf 1_\Delta$, $\tilde h$ a Borel version of
$h$).  Then:
\begin{enumerate}
\item[\rm(i)] $\Cov_{h_\alpha}(Q_1,U_{\Delta,n})=\int_{l-1/2}^{l+1/2}\Cov_{h_\alpha}(Q_1,U_{\Delta,n}(t))\,dt$.
\item[\rm(ii)] For every $t\in[l-\frac12,l+\frac12)$ the pair $(Q_1,U_{\Delta,n}(t))$ consists
of two localized monomials of type (1.1.5), each with all its factors at one time ($s$, resp.\
$t$), with localization squares in $R_s(I)$, resp.\ equal to $\Delta$, separated by the strip
$P_l$ of Lemma~N.38(ii), which is parallel to the time-zero line and has width $w_\Delta\ge1$.
\item[\rm(iii)] Source Fact~2.2 applied to each pair of (ii), with (i) and
Lemma~\ref{lem:slice}(ii), gives
$|\Cov_{h_\alpha}(Q_1,U_{\Delta,n})|\le C_7e^{-m_1w_\Delta}\GJS{Q_1}\GJS{U_{\Delta,n}}$, the bound
that Step~3 of the proof of Theorem~D.3 takes from Lemma~D.2(a) for this pair.
\item[\rm(iv)] Every pair of (ii) is reached by the printed proof directly, on R1 and on R2.
\end{enumerate}
The same holds, with $Q_1=Q_{n,n',k,k'}$ (two factors, both at the time $0$), for the class (P)
applications of Theorem~D.4.
\end{tprop}

\emph{Status.}  (i)--(iii): proved, conditional on A8(iv) at its printed scope for pairs of
one-time monomials across parallel strips, on A3, and on the reading of A8 Scope (d)
(Remark~N.2(iv)) that the constructed objects, here $Q_1$ and the slices $U_{\Delta,n}(t)$, are
GJS's monomials; the same conditions as for Theorem~\ref{thm:TD4} and the record's
Theorem~D.3.  (iv): proved, as a statement about the scope of the source's printed proof.

\begin{proof}
(i) The map $V\mapsto\Cov_{h_\alpha}(Q_1,V)$ is a bounded linear functional on $L^2(\dmu)$:
$|\int Q_1V\,dq_{h_\alpha}|\le Z(h_\alpha)^{-1}\|V\|_2\|Q_1\|_4\|e^{-\Uc_{h_\alpha}}\|_4$ and
$|\int V\,dq_{h_\alpha}|\le Z(h_\alpha)^{-1}\|V\|_2\|e^{-\Uc_{h_\alpha}}\|_2$, finite by
Assumption~A3 ($Z(h_\alpha)>0$, $e^{-4\Uc_{h_\alpha}}\in L^1(\dmu)$) and Lemma~N.28(b)
($Q_1\in L^4(\dmu)$).  $U_{\Delta,n}$ (Definition~N.17) is the localized monomial of type (1.1.6)
with $r=1$, degree $n$ and kernel $\lambda a_nh\mathbf 1_\Delta$ (Lemma~N.18(i)), the
$L^2(\dmu)$-limit of Definition~N.26(a) (Step~3 of the proof of Lemma~D.2), and the kernel is
bounded and vanishes outside $\Delta$.  Lemma~\ref{lem:slice}(i) gives (i).

(ii) $Q_1$ is as stated by Lemma~N.36(ii),(v).  By Lemma~\ref{lem:slice}(ii), $U_{\Delta,n}(t)$ is
a monomial of type (1.1.5) with one factor at the time $t$, localized in $\Delta$.  $P_l$
separates $\Delta_{l,k}$ from the whole row $l_s$ (Lemma~N.38(ii)), hence the localization
squares of the two monomials, which are disjoint because the width is positive.

(iii) For each $t$, Source Fact~2.2 applies: two localized monomials with disjoint localization
squares, individual degrees $1<\bar n$ and $n<\bar n$, separated by a strip of width $w_\Delta$
(so the widest strip has width $\ge w_\Delta$ and $e^{-m_1d}\le e^{-m_1w_\Delta}$, as in Step~4 of
Lemma~D.2), the cutoff $h_\alpha\in\CE$ (Lemma~N.5(ii)), and the printed hypotheses on $m_0$,
$K_1$, $K_2$, $\lambda$ (the standing assumption).  It gives
$|\Cov_{h_\alpha}(Q_1,U_{\Delta,n}(t))|\le C_7e^{-m_1w_\Delta}\GJS{Q_1}\GJS{U_{\Delta,n}(t)}$.
Integrate over $t$ with (i), and use
$\int\GJS{U_{\Delta,n}(t)}dt\le\GJS{U_{\Delta,n}}$ (Lemma~\ref{lem:slice}(ii)).

(iv) Both monomials of the pair are of type (1.1.5), each with all its factors at one time, and
the strip is parallel.  After a time translation (Lemma~N.13 and Remark~N.2(ii): the pieces stay
sharp-time pieces in the rows of their new times, with unchanged norms), and relabelling
$Q_1\leftrightarrow Q_2$ if $l<l_s$ ((1.1.10) is symmetric in $Q_1$ and $Q_2$), the earlier
monomial is localized at times $\le0$ and the later at times $\ge w_\Delta$.  They are admissible
as $A$ and $B_T$ in Theorem~1.1.11 on R2, hence also on R1, and $B_T$ is admissible in
Theorem~1.1.8.  For
Theorem~D.4's class (P) the same proof applies with $Q_1=Q_{n,n',k,k'}$, whose two factors both
sit at the time $0$ (Lemma~N.43(ii)).
\end{proof}

So the slicing of $U_{\Delta,n}$ uses no multi-time reading: its slices have one factor
($r=1$, Definition~N.17), and $Q_1$ has all its factors at one time.  What GJS do not print is
the Fubini step (i), (iii).  The finding N3 is confirmed.  Proposed wording for Remark~N.39, in
place of the clause after the semicolon: ``the printed proof reaches the pair after
$U_{\Delta,n}$ is sliced in time and the bound is integrated over the slices, a Fubini step that
GJS do not print (\path{ugjs/route_t/route_t.tex}, Proposition~T.19); every pair of slices consists of
two monomials of type (1.1.5), each with all its factors at one time, so the pair is reached on
either reading of the equal-times sentence of p.~596, and no multi-time reading is needed.''
Consequence: the uniform gap (Theorem~7.8), which applies Theorem~1.1.7 only through Lemma~7.4
(P1) and Theorem~D.3 (P2), rests only on applications the printed proof reaches.

\section{Status}\label{sec:status}

\subsection{The statements of this note}

\begin{center}\small
\begin{tabular}{@{}p{0.22\textwidth}p{0.72\textwidth}@{}}
\toprule
statement & status\\
\midrule
Definition T.1 & definition\\
Theorem T.2 & proved, conditional on: A1--A8 of the record, with A8(iv) used only in P1--P4,
A8(iv)-T not used, and A2 used at two times placed asymmetrically in a slab (as formula (3) of
the record's Lemma~4.1 uses it at one); the record's declared readings (A1, A3, A6 Scope; A8
Scope (d)) and the open items of \path{part_i/C5_reading_record.md}; the statements of the record
listed in \S\ref{sec:assembly}(c) and in the rows of this table\\
Lemmas T.3, T.4, T.5 & proved (spectral theorem, Fubini, dominated convergence; no import)\\
Lemma T.6 & proved, from Step~2 of the proof of Lemma~9.3 (Proposition~2.7, Locality of \S2.1),
(G$_\gamma$), and Lemmas T.3--T.5\\
Lemma T.7 & proved, from Proposition~8.2(ii) (first display of the proof of Lemma~10.1), the
first assertion of Lemma~10.2, Proposition~2.6, Lemma~9.1(ii), (G$_\gamma$)\\
Lemma T.8, Proposition T.9, Lemma T.10 & proved, from Lemmas T.3, T.6, T.7, Proposition~2.7,
Lemma~3.1, Proposition~8.2(i), Lemma~9.1(i); $\Mt(\lambda_0)$ enters as a number, finite by
Theorem~T.16\\
Lemma T.11 & proved (elementary)\\
Lemma T.12 & proved, from Definition~N.26 and Lemmas N.24, N.25, N.28(a) of the record, Fubini,
Pettis's measurability theorem and the Bochner integral\\
Lemma T.13 & proved, from Lemma T.12 and the record's Lemmas N.20(iv), N.21(ii) (its argument),
N.23(b), N.23$'$(ii)--(iii), N.28(c), N.28(f)(2), and the isometry $J_\tau$ (proof of N.22)\\
Theorem T.14 & proved, conditional on A8(iv) at its printed scope in the cases P3--P4 and on the
reading of A8 Scope (d) (Remark~N.2(iv)); from Lemmas D.1, D.2, C.2, Proposition~N.15, Lemmas
N.5, N.11, N.18, N.32, A3, Lemma~T.13.  Reached by the printed proof of Theorem~1.1.7 on R1
(P4: granted the rotation presupposition), not on R2 (Proposition~T.17(c))\\
Lemma T.15 & proved, from A2 (its multi-time formula for bounded functions on finite slabs, at two
times placed asymmetrically in the slab, derived in Step~1 of (ii) from the bounded case of
Simon's Theorem~V.15 as printed, by approximation), A1(c), A1(d), A3, A5(i), Lemma~N.22 (the
definition of $J_\tau$, its density extension, its isometry), Lemma~N.32, Lemma~3.1(ii),
Lemma~4.4, Lemma~T.13, and the truncation of Lemma~4.3(ii); Simon's Remark on p.~164 is not
used\\
Theorem T.16 & proved, from Lemmas T.11, T.15 and Theorem T.14 (hence conditional as T.14)\\
Proposition T.17 & as a statement about the scope of the source's printed proof (notion of
\S\ref{sec:input}): proved for (a), (b), (d), for P3 on R1, and for the assertions on R2 in (c)
except the reach of P4's degree-$0$ pieces; proved, conditional on the rotation presupposition,
for P4 on R1 and for the reach of P4's degree-$0$ pieces\\
Remark T.18 & its two ingredients are proved: on R2 a transverse pair is reached only if each
monomial has at most one factor of positive degree, and the comparison around a box of fixed shape
needs transverse strips (Lemma~N.7(iv)).  Its conclusion about routes of that type is an argued
observation: the class of routes is not defined\\
Proposition T.19 & (i)--(iii): proved, conditional on A8(iv) at its printed scope (pairs of
one-time monomials across parallel strips), A3 and the reading of A8 Scope (d); (iv): proved, as
a statement about the source's printed proof\\
\bottomrule
\end{tabular}
\end{center}

\subsection{What Theorem 12.2 would rest on if the note were integrated}

\emph{Items (1)--(4) and Corollary~12.3.}  Proved, with the rate function $\Psit_{m_1}$ in
item (1), conditional on the following:
\begin{itemize}
\item A1--A8, with A8(iv) used in the cases P1--P4 only (A8(iv)-T not used), and with A2 used at
two times placed asymmetrically in a slab (Lemma~T.15);
\item the declared readings (A1, A3, A6 Scope; A8 Scope (d); the C5 items);
\item the record's statements listed above.
\end{itemize}
Under the standard ``only cases the printed proof reaches'', they are conditional in addition on
reading R1 of p.~596 (for P3--P4) and, for P4, on the rotation presupposition.  The question of
the constant taken from Theorem~1.1.8 at the same $h$ is common to every case, as in the record.

\emph{Item (5), the gap clause} (Theorem~7.8): unchanged.  It uses P1 and P2 only, which the
printed proof reaches on either reading (Proposition~T.19).  \emph{Item (5), the moment clause}
(Theorem~5.1): unchanged.  It still rests on Theorem~D.4, hence on A8(iv)-T.

\emph{Section 13 of the record.}  Theorem~13.2 uses Lemma~9.3 with the bounded $X=B$
(ll.~3207--3221), Theorem~7.8, and items (1), (3), (4) of Theorem~12.2.  Item (1) enters only
through local-norm convergence (l.~3225); item (4) enters at l.~3227 and items (3)--(4) at
l.~3234.  It uses
neither $M$ nor $\Psi_{m_1}$.  After Route T it rests on what those items rest on.
Remark~13.3 is context.

\subsection{What the note does not do}

\begin{itemize}
\item It does not touch the sharp moment clause of item (5), Theorem~5.1, Corollary~5.2 or
Theorem~D.4; they still rest on A8(iv)-T.
\item It does not prove A8(iv)-T, P3 or P4, does not decide between R1 and R2, does not prove the
rotation presupposition, and verifies no proof of GJS.
\item It does not re-audit the record's statements it cites (Appendix~C, Lemmas D.1, D.2,
N.5--N.43); it uses them as the record states them.
\item It does not evaluate $S_2(B_{3\times5})$; $\Mt_*$ and $\Psit_{m_1}$ contain it through
$M_B$.
\item It is not integrated into \path{part_i/main.tex}.  Integration would touch the following
places of the record (line numbers at 49f2d23).
\begin{itemize}
\item The imports paragraph of Theorem~12.2 (ll.~3083--3099).  It says the theorem uses
A8(iv)-T ``for the sharp-time observable of Theorem~5.1'' and that ``The constant $M$ of
item (5), and through it the rate $\Psi_{m_1}$ of item (1)'' contain $S_2$.  Also the first
sentences of its proof (ll.~3104--3106).
\item A8, item (iv-T), ``Used in'' (ll.~1166--1169), and A8 Scope (b) (ll.~1369--1377).
\item A2: its parenthetical on the uses of the multi-time formulas (ll.~835--838, ``used in \S4
in the two- and one-time forms'') and its ``Used in'' line (l.~848), which would add the
asymmetric two-time use of Lemma~T.15(ii).
\item The ``Integration of Appendix~D'' paragraph of the front matter (ll.~120--124).
\item The introduction: the rate $\Psi_{m_1}$ (ll.~237--241 and 351); the strategy table,
step (i) ``incl.\ A8(iv)-T'' (l.~251) and step (iv) ``(G)+(M$_2$)'' (l.~257); the dependency
display (l.~280); Pillar~1 (ll.~356--368).
\item The notation table: the rows $M$ (l.~613), $\Psi_\gamma$ (l.~618) and $d_0$ (l.~619).
\item Remark~3.3 (ll.~1656--1657 and 1698--1701), which names $M$ and $\Psi_{m_{1,*}}$.
\item Lemma~10.2, split into its decomposition and its bound.  Lemmas~10.3--10.5 and
Proposition~10.4 replaced, and $\Psi_{m_1}$ redefined at the head of \S11 (l.~2921).
\item Appendix~A: the rows $d_0$, shell count, $\Psi_\gamma$ and its tail (ll.~3369--3372), and
l.~3378.
\item Proposition~C.7 (ll.~4087--4101), the inventory of every application of Source Fact~2.2,
which must add P3--P4.
\item The Discussion (l.~3291), which calls $\Psi_{m_1}$ ``super-polynomial but sub-exponential''.
$\Psit_{m_1}$ contains the same $T_{m_1}$ terms; that characterization was not checked for
either function.
\item The new statements T.3--T.16, and the wording of Remark~N.39 above.
\end{itemize}
\item A point for the record, not for this note.  The approximation in Step~1 of the proof of
Lemma~T.15(ii) applies equally to the record's own Lemma~4.1.  A2 is stated ``For all bounded
time-zero functions'', while its printed sources give the interior functions as functions of
finitely many fields (Simon's Theorem~V.15; GRS's Theorem~II.16).  For other bounded functions,
Lemma~4.1 rests on Simon's Remark on p.~164 or on this approximation.
\end{itemize}

\subsection{The steps an auditor of this revision should check first}

\begin{enumerate}
\item Lemma~T.15(ii), rewritten twice.  Check that A2's ``analogous multi-time formulas on finite
slabs'' cover bounded functions at two asymmetric times.  Check the approximation of Step~1 by
bounded cylinder functions, which rests that formula on the bounded case of Simon's
Theorem~V.15 as printed.  Check the identity $J_\tau(c_\ell\circ F)=c_\ell\circ J_\tau F$ of
Step~0.  Check the two dominated convergences of Step~2.
\item Proposition~T.17(c), now written out for P3: the $L^4$-Bochner representation of the slices
and the two-step use of the bounded covariance functional.
\item The account of the printed evidence for R1 in \S\ref{sec:readings}, in particular the
reading of p.~628.
\item Remark~T.18, now stated as an argued observation.
\end{enumerate}

\appendix
\section{Numerical checks (numerical; nothing is deduced from them)}\label{app:numerics}

The script \path{check_route_t.py}, kept beside this note, was run on 23 November 2026 with
the system \texttt{python3} (Python 3.14.7, numpy 2.5.2, scipy 1.18.1, mpmath 1.4.1), seed
$20261123$; its output, \path{check_route_t.output.txt}, ends with ``ALL CHECKS PASSED''.  It checks:
the closed form of $k$ against two quadratures at nine values of $E\in[10^{-6},100]$ (40-digit
arithmetic, relative tolerance $10^{-12}$; maximal error $9.2\cdot10^{-30}$); the inequality of
Lemma~T.11 on a grid of $4\cdot10^5$ points $E\in[0,10^4]$ for $\lambda_0\in\{1,1.0001,1.5,3,10\}$
(minimal margin $3.5\cdot10^{-4}$) and the monotonicity of $k$; the transform of
$f_{a,\varepsilon}$ (Lemma~T.5, Step~1) at four parameter sets (error $<10^{-9}$); the
principal-value limit of Step~3 at three parameter sets, for $\varepsilon=10^{-2},10^{-3},10^{-4}$
(errors decreasing like $\varepsilon$, the smallest $\le1.1\cdot10^{-4}$) and the bounds on $J$;
the Gaussian tail $1-\Phi(x)\le e^{-x^2/2}/2$ on $[0,40]$; $\int_0^x\sinh u\,u^{-1}du\le\sinh x$;
the resolvent transform and the identity of Lemma~T.7, Step~3, for a random discrete measure; the
discrepancy formula of Lemma~T.4 for two functions that agree on $(-2,2)$; the separator tail of
Step~2; the count of integers in a half-open interval of length $3$; the monotonicity of $\Xi$;
$\int\|w(t,\cdot)\|dt\le\|w\|$ for random step kernels; and, for the free field with $n=1$,
$m_0=1$, $h=\mathbf 1_{[0,2]}$, the inequality $\E[\Vb(h)^2]\le\E[V(h)^2]$ ($0.478\le0.709$).  The
tolerances are printed in the output.

Two checks were added with the revision, and the script was rerun on 23 November 2026 (23 lines
PASS, ``ALL CHECKS PASSED'').
\begin{itemize}
\item \emph{The truncation of Lemma~T.15(ii), Step~2,} in a Gaussian toy: $X_i=Y_i^2-1$ for
standard normal $Y_1,Y_2$ with correlation $\rho$.  The quantity $\E[X_1^{[\ell]}X_2^{[\ell]}]$
is computed by the Mehler expansion in normalized Hermite polynomials, 120 terms, and compared
with $\E[X_1X_2]=2\rho^2$.  For $\rho\in\{0.3,0.7,0.9\}$ and $\ell=0.5,2,8,32,64$ the errors
decrease, to at most $2.2\cdot10^{-13}$ at $\ell=64$.  The first run of this check stopped at
$\ell=32$ with an a priori tolerance of $10^{-6}$.  It printed FAIL because the error at
$\ell=32$, $\rho=0.9$, is $1.03\cdot10^{-6}$: a true truncation error, not a numerical one.
The level $\ell=64$ was then added.
\item \emph{The free case of Lemma~T.15(ii)} ($\lambda=0$).  Here the formula reads
\[
 \E_{\mu_0}[\Phi(\tau,f)\Phi(\tau',f)]=\ip{\varphi_0(f)}{e^{-(\tau'-\tau)H_0}\varphi_0(f)} .
\]
The
position-space value $\iint f(x)f(y)\,K_0(m_0|(u,x-y)|)/(2\pi)\,dx\,dy$ is compared with the
momentum-space value $(2\pi)^{-1}\int|\hat f(k)|^2e^{-u\mu(k)}/(2\mu(k))\,dk$, for
$f=\mathbf 1_{[0,1]}$, $m_0=1$, $u=\tau'-\tau\in\{0.3,1,2.5\}$.  The relative difference is
$1.7\cdot10^{-15}$ (tolerance $10^{-8}$).
\end{itemize}

\end{document}
